% =============================================================================
% Project: 302006 - OOD-robustness benchmark of single-cell foundation models
% File: main.tex (P2 rewrite, assembled 2026-09-11 from revision/drafts/ +
%   P0_onepager.md round-17 final numbers; supersedes the 2026-08-26 version,
%   which remains in git history)
% Title/framing: "benchmarking review of five published models" (user decision
%   2026-09-09, BIB scope). All statistics trace to P0_onepager.md or the
%   results CSVs cited per section.
% Compile: pdflatex main (x2). Supplementary cross-refs resolve after the P2
%   supplementary.tex pass (xr).
% =============================================================================

\documentclass[11pt,a4paper]{article}

% Packages
\usepackage[utf8]{inputenc}
\usepackage[T1]{fontenc}
\usepackage{amsmath,amssymb,amsfonts}
\usepackage{graphicx}
\usepackage{booktabs}
\usepackage{multirow}
\usepackage{array}
\usepackage{xcolor}
\usepackage{hyperref}
\usepackage[margin=1in]{geometry}
\usepackage[numbers,sort&compress]{natbib}
\usepackage{lineno}
\usepackage{float}
\usepackage{subcaption}
\usepackage{authblk}
\usepackage{xr}

% Resolve \ref to labels defined in supplementary.tex.
% Compile order: pdflatex supplementary -> pdflatex supplementary -> pdflatex main -> pdflatex main
\externaldocument{supplementary}

% Custom commands
\newcommand{\model}[1]{\textsc{#1}}
\newcommand{\auc}{\text{AUC}}
\newcommand{\ci}[2]{[#1, #2]}

% Line numbers for review
\linenumbers

\title{Stress-testing single-cell foundation models on underrepresented
platelets: a patient-level mortality benchmark}

\author[1]{Xinru Qiu}
\author[1]{Adam Godzik}
\affil[1]{Division of Biomedical Sciences, School of Medicine, University of California, Riverside, CA, USA}

\date{}

% Allow TeX to loosen inter-word spacing rather than overrun the margin
\setlength{\emergencystretch}{3em}

\begin{document}

\maketitle

% ============================================================================
% ABSTRACT (P0_onepager.md Sec.4 final, 308 words)
% ============================================================================
\begin{abstract}
We benchmarked five single-cell foundation models (\model{STATE}, \model{UCE},
\model{Geneformer}, \model{scPRINT-2}, \model{scGPT}) on human platelets, a
minority cell type in the public atlases these models draw their pretraining
corpora from. We used in-hospital and 28-day mortality as a patient-level task
in 240 COVID-19 and 91 sepsis patients. All models received identical raw
input and were evaluated under repeated patient-level cross-validation, every
prediction aggregated to the patient before scoring, against a from-scratch
PCA + logistic-regression baseline, a null model, and an in-distribution
monocyte positive control. Mortality was learnable from platelet
transcriptomes (patient-level AUC up to 0.79), but under a pre-specified,
selection-corrected test no foundation model significantly outperformed the
baseline (label-permutation $P = 0.181$; cross-check $P = 0.073$; stable
across all 25 leave-one-fatal-out refits). A separate question --- whether a
model's score adds information beyond the baseline's even without ranking
patients better on its own --- was examined with a nested-model comparison; no
model's increment survived the same multiple-comparison correction used
throughout (Benjamini--Hochberg, 0/5 in every comparator configuration
tested). Within sepsis, no model exceeded chance at late sampling windows (17
events, resolving only AUC $\gtrsim 0.71$), consistent with published
platelet-outcome signals being reported only in early disease windows.
Cross-disease transfer was rank-indeterminate in the direction where signal
exists and uninformative in the direction where it does not; the adaptation
curve has the same resolution limit. The one result that survived
selection-corrected and multiple-comparison testing was that two models
(\model{STATE}, \model{Geneformer}, statistically indistinguishable from each
other) separated from a third (\model{UCE}), the weakest method throughout ---
the only one that failed to clear the null model built from cell counts and
disease identity. At realistic clinical event counts, pretraining bought no
detectable advantage over a strong from-scratch baseline on this
underrepresented cell type; the benchmark quantifies the resolution achieved
($\sim$0.10 AUC at 25 events) and the scale needed to detect smaller effects.

\textbf{Keywords:} single-cell foundation models; benchmarking; patient-level
evaluation; platelet transcriptomics; mortality prediction; statistical
resolution
\end{abstract}

% ============================================================================
% KEY POINTS (BIB-required; P0_onepager.md Sec.7 final wording)
% ============================================================================
\section*{Key Points}
\begin{itemize}
    \item Five published single-cell foundation models were benchmarked on
    platelet mortality prediction; none significantly outperformed a
    from-scratch PCA+logistic-regression baseline under a pre-specified,
    selection-corrected test (permutation $P = 0.181$).
    \item A pretraining-corpus-composition hypothesis for the one confirmatory
    finding was overturned by the corpus measurement itself --- the three
    models whose corpora can be measured show no ordering by platelet content.
    \item The one result that survives correction is not that two models have
    a pretraining advantage --- it is that a third (\model{UCE}) is the
    weakest performer throughout, the only one that fails to clear the
    count+disease null model.
    \item The same null held on an in-distribution positive control
    (monocytes), showing platelet underrepresentation is not the explanation
    --- the benchmark's pipeline finds no foundation-model advantage anywhere
    it looked.
    \item The paper's contribution is a resolution-quantified benchmarking
    protocol ($\sim$0.10 AUC resolvable at 25 events), not a leaderboard ---
    a distinction that matters because comparable patient-level scFM
    evaluations in the literature report no correction for the comparisons
    they make, and none reports event counts at the granularity of those
    comparisons.
\end{itemize}

\newpage

% ============================================================================
% INTRODUCTION
% ============================================================================
\section{Introduction}

Single-cell foundation models (scFMs) promise general-purpose representations
of cellular state learned by self-supervision on tens of millions of cells
from public reference atlases \citep{cui2024scgpt, theodoris2023geneformer,
regev2017hca, tabula2022sapiens, czi2025cellxgene}, importing into biology the
paradigm that transformed language and vision \citep{bommasani2021foundation,
moor2023generalist, bunne2024virtualcell}. The family has grown rapidly and
now spans masked-language, rank-value, generative, and ontology-aware
objectives \citep{cui2024scgpt, theodoris2023geneformer, hao2024scfoundation,
wen2024cellplm, rosen2026universal, kalfon2025scprint, adduri2025state,
yang2022scbert, chen2025genept, tejada2025nicheformer},
reviewed comprehensively by \citet{szalata2024transformers} and, most
recently, by a field-wide landscape analysis \citep{jiang2026landscape}. What
these models are good for, however, depends on how they are evaluated, and the
evaluation literature splits sharply in two. On cell-level tasks --- cell-type
annotation, clustering, batch integration, perturbation response ---
benchmarking is mature and has already produced a sobering pattern: zero-shot
foundation models are matched or beaten by simpler methods such as scVI,
highly-variable-gene selection, or mean expression
\citep{kedzierska2025zeroshot, ahlmanneltze2025linear, bendidi2024pca,
luecken2019bestpractices, luecken2022benchmarking, luecken2025openproblems,
csendes2025benchmark, wenteler2025perteval, boiarsky2023deepdive,
denadel2026scaling, liu2026evaluating}. On
patient-level clinical-outcome prediction, the setting in which scFM
representations would have to justify their scale for precision medicine
\citep{jiang2026landscape}, systematic evaluation is essentially absent. The
few studies we could identify (Table~\ref{tab:review_outcomes}
\citep{egsp2026survival, empiricalscfm2025cancer, tissueformer2025covid})
share three gaps: none reports a correction for the comparisons it makes (winner
selection, multiplicity), none discloses event counts at the granularity of
its own comparisons, and event counts are what determine whether any two
scores are distinguishable at all.

The present study is a benchmarking review of five published single-cell
foundation models --- \model{STATE} \citep{adduri2025state}, \model{UCE}
\citep{rosen2026universal}, \model{Geneformer}
\citep{theodoris2023geneformer}, \model{scPRINT-2} \citep{kalfon2025scprint},
and \model{scGPT} \citep{cui2024scgpt} --- under one standardized,
patient-level protocol, rather than a new-method paper. As the test case we
pair a clinically decisive endpoint, mortality, with a cell type that the
pretraining corpora barely contain. Platelets are anucleate cell fragments
carrying megakaryocyte-derived mRNA
\citep{rowley2011genomewide, bray2013complex, zaid2020platelets}; in the
three corpora whose composition
can be numerically bounded against public records they account for
approximately 0.12\% of cells, \model{STATE}'s corpus is bounded lower still
(below 0.05\%) through its public component documentation, and only
\model{Geneformer} publishes no cell-type-level corpus composition at all.
They are also mechanistically proximal to fatal inflammatory disease:
platelet-driven immunothrombosis is implicated in COVID-19 and sepsis
mortality \citep{manne2020platelet, barrett2021platelets, hottz2020platelet,
ackermann2020pulmonary, nicolai2020immunothrombotic, levi2020coagulation,
shappell2023sars},
and outcome-predictive platelet transcriptomic signals have been reported in
early sepsis windows
\citep{qiu2021sepsis, middleton2019sepsis, middleton2020nets, qiu2024deciphering}. We assemble
240 COVID-19 patients (25 fatal) and 91 sepsis patients (17 fatal) with
per-patient outcome labels from public platelet scRNA-seq cohorts, and we ask
four questions. Does any foundation model beat a strong from-scratch baseline
(PCA + logistic regression on the same cells) at patient-level mortality
prediction, once winner selection is accounted for? Is the signal real at
all, or reducible to cell counts, library quality, or cohort membership? Does
it transfer across diseases and cohorts? And when the models are mapped back
to genes, do they even use the same biology
(Figure~\ref{fig:overview})?

The answers are uniformly sobering and quantitatively bounded. Mortality is
learnable from platelet transcriptomes beyond every artifact control we could
build, but under a pre-specified, winner-corrected test no foundation model
outperformed the from-scratch baseline (permutation $P = 0.181$), and a
nested complementarity analysis returned the same null under
multiple-comparison correction. The same null returned within sepsis, in both
cross-disease directions, and on an in-distribution monocyte positive control,
and the one confirmatory separation that survives correction is a statement
about the weakest model, not the strongest. Because every null is reported
with the resolution of its design ($\sim$0.10 AUC at 25 fatal events;
detecting 0.03--0.05 AUC effects requires on the order of $10^3$ patients at
this event rate), the benchmark doubles as a plug-in protocol: any future
model can be dropped into the identical pipeline and read against the same
baselines, and any future claim of a pretraining advantage on this task now
has a measured bar to clear.

% Table: review content (P0_onepager.md Sec.5; search strategy verbatim)
\begin{table}[t]
\centering
\caption{Patient-level clinical-outcome evaluations of single-cell foundation
models identified by our search (PubMed via Google
\texttt{site:pubmed.ncbi.nlm.nih.gov} + bioRxiv/medRxiv; query
``\texttt{single-cell foundation model patient-level clinical outcome
benchmark survival mortality 2025 2026}'', run 2026-09-09; 4 unique records
recalled, 2 included, 2 excluded as not scFM-based). TissueFormer was
identified in an earlier literature pass and confirmed retrievable by a
targeted search rather than the general query. None of the three reports a
selection or multiplicity correction for its comparisons, and only the first
discloses pooled event counts.}
\label{tab:review_outcomes}
\small
\resizebox{\textwidth}{!}{%
\begin{tabular}{p{4.2cm}p{2.6cm}p{3.4cm}p{2.2cm}p{1.9cm}p{1.6cm}}
\toprule
\textbf{Study (venue)} & \textbf{scFM(s)} & \textbf{Task / outcome} & \textbf{From-scratch baseline} & \textbf{Patient-level aggregation} & \textbf{Correction} \\
\midrule
EGSP, Bioinformatics Advances 2026 (peer-reviewed; PMC13032892) & scFoundation (embedding transfer to bulk) & pan-cancer overall survival (TCGA, 25 types; 8{,}776 samples) & yes (Ridge/Lasso-Cox) & yes & none reported \\
Empirical Evaluation of scFMs for Predicting Cancer Outcomes, bioRxiv 2025.10.31.685892 (preprint) & 9 models & 6 cancer clinical/biological outcome tasks (1{,}170 experiments) & yes (HVG, PCA, scVI) & yes (3 schemes) & none reported \\
TissueFormer, bioRxiv 2025.08.17.670735 (preprint) & murine-adapted Geneformer & COVID-19 severity, 3-class (376 donors, 3 cohorts) & yes (random-init) & yes & none reported \\
\bottomrule
\end{tabular}%
}
\end{table}

\begin{figure}[t]
\centering
\includegraphics[width=\textwidth]{figures/Figure1_overview.pdf}
\caption{\textbf{Study design.} \textbf{(A)} Cohorts: platelet single-cell
transcriptomes from 240 COVID-19 patients (25 fatal), 103 sepsis patients in
the dataset (19 fatal; the seven GSE151263 patients, who lack per-patient
outcome records, are excluded from all mortality analyses), and 70 healthy
donors (reference only); a monocyte in-distribution positive control covers
148 patients (19 fatal). Multiple draws per patient are merged to the patient
level (61{,}759 platelet transcriptomes in the full COVID-19 and sepsis
dataset; 61{,}283 across the mortality-analysis pools after the GSE151263
exclusion). \textbf{(B)} Unified input: every model receives the
identical raw AnnData object and applies its own native preprocessing; the
five foundation models are benchmarked frozen, against a from-scratch
PCA-200 + logistic-regression baseline and four control families.
\textbf{(C)} Patient-level evaluation: repeated cross-validation (10 seeds
$\times$ 5 folds) with patients grouped so that no patient's cells or draws
span folds; cell-level predictions are aggregated to one score per patient;
the primary metric is patient-level out-of-fold ROC-AUC, under a three-tier
inference hierarchy (pre-specified primary, confirmatory secondary,
exploratory). \textbf{(D)} The four analysis arms and their headline
results. Figure created with BioRender.com.}
\label{fig:overview}
\end{figure}

% Methods (synced to P0 round-17 state, 2026-09-09)
\section{Methods}
\label{sec:methods}

\subsection{Cohort and endpoint definitions}
\label{ssec:methods_cohort}

We analyzed human platelet single-cell transcriptomes from public cohorts
(Table~\ref{tab:cohort}): 240 COVID-19 patients (25 fatal), 103 sepsis
patients in the dataset (19 fatal), and 70 healthy donors retained as a
reference group only. Multiple draws of the same patient (longitudinal samples
or replicate libraries) were merged into a single patient-level unit
(\texttt{patient\_id}) before any analysis, so that all of a patient's cells
and draws share one cross-validation group and one aggregated score.
GSE151263 (7 patients) does not report per-patient outcomes in the source
record and is excluded from all mortality analyses; it is retained in the
dataset for completeness. This eligibility decision follows from source-record
availability and was fixed at dataset assembly, independently of any
outcome-associated result. No reported table or figure contains GSE151263
patients; the cohort appears only in the dataset description
(Table~\ref{tab:cohort}). In arithmetic terms: of the 103 sepsis patients in
the dataset, 7 enter no outcome-labeled analysis, 91 (GSE279452) form the
main sepsis arm, and all 96 outcome-labeled patients (91 plus 5 from
GSE167363) form the reverse-transfer training pool.
The primary endpoint was mortality: in-hospital
death for the COVID-19 cohorts and GSE167363, and 28-day death for GSE279452;
the difference is stated wherever the GSE279452 arm is used. Survivors
include patients discharged alive at any severity grade (severe cases are
ICU-discharged survivors, labeled as the negative class). Outcome labels were
verified against the source records for every fatal patient
(Supplementary Note~\ref{sn:label_provenance}).

\subsection{Foundation models and unified input}
\label{ssec:methods_models}

We evaluated five single-cell foundation models (scFMs) chosen to span the
main representation families (Table~\ref{tab:models}): \model{STATE},
\model{UCE}, \model{Geneformer}, \model{scPRINT-2}, and \model{scGPT}. To
remove input preprocessing as a confounder, every model received the identical
raw AnnData object \citep{wolf2018scanpy} and applied its own native
preprocessing pipeline as recommended by its authors; embeddings were then used frozen, with no
fine-tuning, generative, or task-specific capabilities of any model. For
\model{STATE}, a two-component system, we evaluate the state-embedding
component (SE-600M): the state-transition module is trained for
perturbation-response prediction and has no natural mapping to a
patient-level discrimination task, and SE is the system's representation
component. Platelets are anucleate cell fragments carrying
megakaryocyte-derived mRNA; in the three pretraining corpora whose
composition can be numerically bounded they account for approximately
0.12\% of cells, and \model{STATE}'s corpus is bounded lower still (below
0.05\%) through its public component documentation, while only
\model{Geneformer} publishes no cell-type-level corpus composition at all
(Table~\ref{tab:models} and Supplementary Note~\ref{sn:corpus_platelet}).
They are therefore an underrepresented substrate in every corpus that can
be checked --- documented as absent from some, bounded near one per mille
in the rest.

\subsection{Evaluation protocol}
\label{ssec:methods_cv}

The primary evaluation unit is the patient. Classifiers were trained and
evaluated under repeated patient-level cross-validation: 10 repeats of 5-fold
\texttt{StratifiedGroupKFold} (\texttt{shuffle=True}, seeds 42--51) with
\texttt{patient\_id} as the group variable, so that no patient's cells or
draws ever span training and test folds
\citep{varma2006bias, kapoor2023leakage, whalen2022pitfalls,
mcdermott2021reproducibility, koh2021wilds}. Within each repeat, out-of-fold cell
predictions were averaged per patient (mean; median and trimmed-mean as
sensitivity analyses, Supplementary Note~\ref{sn:aggregation}) to yield one
score per patient. The primary metric is patient-level out-of-fold ROC-AUC
(patient-level AUC) computed as the repeat-mean, reported with the standard
deviation across repeats (repeat-sd) as a measure of split-to-split stability.

We pre-specified a single primary inferential question: does the best
foundation model beat the from-scratch baseline? Because quoting the best of
five correlated candidates inflates the error rate, this contrast was tested
with a winner-corrected max-statistic permutation test: mortality labels were
permuted at the patient level (10{,}000 permutations) on the frozen
seed-ensemble scores, and the familywise $P$ is the fraction of permutations
whose maximum foundation-model-minus-baseline AUC difference reaches the
observed maximum; a recentered bootstrap is reported as a cross-check
(Supplementary Note~\ref{sn:winner_test}).
Every other contrast is reported as
exploratory, with nominal $p$-values from exact DeLong's test \citep{delong1988comparing}
labeled as such
(paired-bootstrap agreement documented in Supplementary
Note~\ref{sn:se_check}), and the stability of the primary result was assessed by
leave-one-fatal-patient-out refits (Supplementary Note~\ref{sn:fragility}).
Uncertainty quantification was matched to each analysis
design: repeated-CV for fold-based within-disease evaluation, patient-clustered
bootstrap (200 iterations) for fixed-split cross-disease transfer, and
training-pool subsampling (10 seeds) for the adaptation curve.
All feature-wise preprocessing was fit inside training folds only:
standardization, PCA, and any gene selection saw no test cells, and the 70
healthy donors entered no model fit. The primary classifier was logistic
regression ($C=1.0$, lbfgs, balanced class weights) and the secondary a random
forest (100 trees, balanced class weights), both as implemented in
scikit-learn \citep{pedregosa2011scikit}.
Two safeguards apply to every inferential construction. Any
selection-corrected $P$ must be at least the smallest uncorrected $p$ in its
comparison family, and any nested out-of-fold pooling harness must pass a
unit test, in which a known frozen score refit through the harness must
reproduce its own out-of-fold AUC within 0.02--0.04, before it may produce
significance conclusions (Supplementary Note~\ref{sn:nested_comparison}).

\subsection{Baselines and controls}
\label{ssec:methods_baselines}

Six controls ran under the identical protocol. (i) A from-scratch
PCA + logistic-regression baseline on raw counts: PCA dimensionality was
selected from \{50, 100, 200, 500\} in a pre-registered sweep with
selection performed inside training folds, and we report the 200-component
setting throughout (PCA-200+LR). (ii) A null model using only each patient's
cell count and disease label; within COVID-19, where disease is constant, it
reduces to a cell-count-only model. (iii) A QC-metrics-only model using
per-patient mean UMI counts, detected genes, percent mitochondrial reads, and
cell count, to test whether predictions could be explained by library-quality
or processing artifacts (Supplementary Note~\ref{sn:qc_baseline}).
(iv) A cohort-only model scoring each patient by their cohort's mortality
rate, estimated within training folds only
(Supplementary Note~\ref{sn:cohort_baseline}). (v) A nested incremental-value
family, in which a covariate baseline (QC metrics, cell count, and cohort
identity; L2-regularized logistic regression at $C = 0.3$) is augmented with
each representation score in turn; its pre-specified primary is the single
PCA-200 nested likelihood-ratio test against an empirical permutation null,
chosen because it involves no model selection
(Supplementary Note~\ref{sn:qc_incremental}), and each method is additionally
compared with the null model directly
(Supplementary Note~\ref{sn:method_vs_null}). (vi) An
in-distribution positive control on monocytes from the same source cohorts
(148 patients, 19 fatal), a cell type the models were pretrained on
(Section~\ref{ssec:r4_monocyte}); monocyte identities came from the published
atlas annotations (Ren, Stephenson) or a marker-based call (GSE167363), and a
patient-matched subset analysis is reported in
Supplementary Note~\ref{sn:monocyte_matched}.

\subsection{Cross-disease transfer and adaptation curve}
\label{ssec:methods_transfer}

Cross-disease transfer used fixed disease splits in both directions
(COVID-19$\rightarrow$sepsis and sepsis$\rightarrow$COVID-19) under the primary
severe-as-survived label definition: the forward direction trained on COVID-19
(240 patients) and tested on the 91 GSE279452 sepsis patients, and the reverse
direction trained on all 96 outcome-labeled sepsis patients and tested on
COVID-19, with patient-level scoring and patient-clustered bootstrap
confidence intervals (200 iterations); a
severe-excluded variant is retained as a sensitivity analysis
(Supplementary Note~\ref{sn:severe}). The adaptation curve measured whether
transfer is recovered with target-disease supervision: COVID-trained
classifiers received increasing numbers of sepsis training cells
($n \in \{0, 10, 50, 100, 500, 1000, \text{all}\}$, 10 seeds per size) and were
scored on a held-out 32-patient sepsis test set, a one-third stratified
holdout of the 96-patient labeled sepsis pool.

\subsection{Metrics and calibration}
\label{ssec:methods_metrics}

The primary metric is patient-level out-of-fold ROC-AUC
(Section~\ref{ssec:methods_cv}). Secondary quantities, always named, are the
cell-weighted out-of-fold AUC, the per-fold mean AUC, the leave-one-cohort-out
(LOBO) cohort AUC, and the out-of-fold nested $\Delta$AUC of the
complementarity family (Supplementary Note~\ref{sn:nested_comparison}); no
unqualified ``AUC'' is used elsewhere. Because cohort is
confounded with disease and sampling site in these data, LOBO is treated as a
secondary robustness axis rather than a co-primary estimand; the fixed-split
cross-disease analysis isolates the transfer question directly
(Section~\ref{ssec:methods_transfer}). A LOBO fold is computable only when the
held-out cohort contains at least one fatal and at least one survivor patient;
cohorts failing this rule (34 patients, all in survivor-only cohorts) remain
in the training pool throughout, so pooled LOBO covers 206 of the 240 patients
with all 25 fatal events retained. A fifth named quantity, the
seed-ensemble out-of-fold AUC (ensemble-OOF AUC), averages each patient's
scores across the ten repeats and scores them once; the primary test, all
paired-bootstrap contrasts, fragility refits, score correlations, and the
consensus predictor are computed on these ensemble scores, while the
repeat-mean and repeat-sd report split-to-split stability of the same
protocol.
Calibration was
assessed on patient-level out-of-fold predictions with the Brier score \citep{brier1950verification} and the
expected calibration error (ECE, 10 equal-width bins)
\citep{guo2017calibration, vancalster2019calibration}. Multiplicity was
controlled within each comparison family with Benjamini--Hochberg FDR
\citep{benjamini1995controlling}, and the
family is stated with every reported $q$.

\subsection{Reproducibility and availability}
\label{ssec:methods_repro}

All analyses ran with fixed seeds (CV seeds 42--51; bootstrap seed 42) and are
fully scripted \citep{fahsbender2025czi}. Code, per-model environments, the patient-level mapping, and
all result tables are available at \url{https://github.com/xqiu625/302006_updategit}. Embeddings were computed
on the UCR HPCC cluster with each model's recommended environment and publicly
released checkpoints (Table~\ref{tab:models}).

\subsection{Masking-based gene attribution}
\label{ssec:methods_masking}

For each model, cells of the canonical within-COVID cohort (240 patients, 25
fatal) were subsampled to at most 60 per patient (10{,}737 cells). Per cell,
the top-256 expressed genes were selected from the intersection of model
vocabularies (deterministic, identical across models sharing the same
universe). One in-sample L2 logistic regression ($C = 0.3$) per model on that
model's own candidate-restricted embeddings defined the mortality direction,
$s = w \cdot z + b$. Each candidate gene was then occluded in turn via the
model's native attention mechanism (attention mask set to 0 for
\model{Geneformer}; physical sequence deletion for \model{STATE}, the only
architecture whose mask flag does not reach attention;
\texttt{src\_key\_padding\_mask} set for \model{scGPT} and \model{UCE}, the
latter with a one-line float-to-boolean cast required for the mask to
function at all; \texttt{mask\_zeros} cleared for \model{scPRINT-2}), and
$\Delta(g, \text{cell}) = s_\text{full} - s_\text{masked}(g)$ was averaged
across cells (minimum 30 occurrences) into a ranked per-gene list per model.
Prerank GSEA \citep{subramanian2005gsea} (gseapy \citep{fang2023gseapy};
MSigDB Hallmark 2020, KEGG 2021, Reactome 2022; 10--500 genes per term;
1{,}000 permutations; seed 20260910) produced per-model pathway tables. The joint-occlusion control masked the top-$N$ attributed
genes jointly ($N \in \{10, 30, 128\}$) versus matched random-$N$ sets (three
draws per cell) over the full cohort, comparing paired per-cell deltas.
% Table 1 — cohort characteristics (counts verified against donor_meta_with_patient_id_v2.csv + Job M)
\begin{table}[t]
\centering
\caption{Source cohorts. Patients are unique patient-level units after merging
multiple draws per patient. ``Fatal'' counts are patients with a fatal outcome
under each cohort's outcome definition. Sampling windows are reported from the
source publications or repository records; ``not reported'' means the source
states no draw timing. GSE151263 is excluded from all mortality analyses
(no per-patient outcome in the source). Healthy donors ($n=70$) are used as a
reference group only and enter no model fit.}
\label{tab:cohort}
\small
\resizebox{\textwidth}{!}{%
\begin{tabular}{lllrrlll}
\toprule
\textbf{Cohort (accession)} & \textbf{Disease} & \textbf{Patients} & \textbf{Fatal} & \textbf{Cells} & \textbf{Sampling window} & \textbf{Platform} & \textbf{Outcome definition} \\
\midrule
Ren 2021 \citep{ren2021covid} (GSE158055) & COVID-19 & 54 & 9 & 8{,}394 & 1--126 d post onset (median 23) & 10x & in-hospital death \\
Stephenson 2021 \citep{stephenson2021singlecell} (E-MTAB-10026) & COVID-19 & 90 & 8 & 11{,}387 & not reported & 10x & death \\
Liu 2021 \citep{liu2021timeresolved} (GSE161918) & COVID-19 & 31 & 4 & 2{,}050 & T0 median 11 d post onset (T0--T3) & 10x & death \\
Schulte-Schrepping 2020 \citep{schulteschrepping2020severe} (EGA) & COVID-19 & 18 & 3 & 616 & 3--20 d (c1), 2--25 d (c2) post onset & 10x / Rhapsody & death \\
Bernardes 2020 \citep{bernardes2020longitudinal} (FastGenomics) & COVID-19 & 13 & 1 & 6{,}030 & d 0/2/7/10/13 and/or discharge & 10x & death \\
Combes 2021 \citep{combes2021global} (GSE163668) & COVID-19 & 13 & 0 & 1{,}985 & screened $\leq$3 d of hospitalization & 10x & (survivors) \\
Wilk 2020 \citep{wilk2020singlecell} (GSE150728) & COVID-19 & 6 & 0 & 180 & 2--16 d post onset & 10x & (survivors) \\
Lee 2020 \citep{lee2020immunophenotyping} (GSE149689) & COVID-19 & 8 & 0 & 944 & not reported & 10x & (survivors) \\
Arunachalam 2020 \citep{arunachalam2020systems} (GSE155673) & COVID-19 & 7 & 0 & 4{,}490 & 2--16 d post onset (median 9) & 10x & (survivors) \\
\midrule
Qiu 2021 \citep{qiu2021sepsis} (GSE167363) & Sepsis & 5 & 2 & 2{,}946 & 0 h and 6 h post sepsis recognition (ICU) & 10x & in-hospital death \\
GSE279452 & Sepsis & 91 & 17 & 22{,}261 & median 16 d post admission (3--121) & 10x & 28-day death \\
Jiang 2020 \citep{jiang2020ards} (GSE151263) & Sepsis/ARDS & 7 & --- & 476 & within 24 h of intubation & 10x & not reported (excluded) \\
\midrule
Healthy donors (multiple) & Healthy & 70 & --- & 3{,}205 & --- & 10x & reference only \\
\bottomrule
\end{tabular}%
}
\end{table}
% Table 2 — foundation models (all facts verified against source papers; see model_architecture_analysis.md)
\begin{table}[t]
\centering
\caption{Foundation models evaluated. All models received the identical raw
AnnData object and applied their own native preprocessing; embeddings were used
frozen with a logistic-regression head. Corpus platelet status summarizes our
verification of each pretraining corpus (Supplementary Note~\ref{sn:corpus_platelet}).}
\label{tab:models}
\small
\resizebox{\textwidth}{!}{%
\begin{tabular}{llllll}
\toprule
\textbf{Model (checkpoint)} & \textbf{Architecture} & \textbf{Training objective} & \textbf{Pretraining corpus} & \textbf{Embedding} & \textbf{Corpus platelet status} \\
\midrule
STATE (SE-600M) & dense bidirectional transformer (16 layers, 600M) & dual-axis expression reconstruction + dataset classification & 167M cells (scBaseCount, CxG, Tahoe-100M) & 2{,}058-dim & absent (scBaseCount annotation space; QC-filtered) \\
UCE (650M) & transformer (33 layers, 650M), ESM2 gene tokens & masked gene modeling (BCE, 20\%) & IMA, 36M cells, 313 datasets & 1{,}280-dim & absent (51-type composition, SI Table 19) \\
Geneformer (V2-316M) & BERT encoder (18 layers, 316M) & rank-value masked-gene modeling & Genecorpus-104M (non-cancer) & 1{,}152-dim & not reported; corpus rich in blood/immune datasets \\
scPRINT-2 (small-v2) & multi-task transformer + ontology classifiers (20M active) & denoising + label prediction + bottleneck learning & 350M cells, 16 organisms & 424-dim & not reported (Census 0.120--0.125\% bound) \\
scGPT (whole-human) & generative transformer (12 layers), generative attention mask & masked expression prediction + GEPC & 33M non-disease cells (441 studies) & 512-dim & not reported (Census $\leq$0.125\% bound) \\
\bottomrule
\end{tabular}%
}
\end{table}

% ============================================================================
% RESULTS
% ============================================================================
\section{Results}
\label{sec:results}

% R1 — within-COVID mortality discrimination (synced to P0 round-17 state, 2026-09-09)
\subsection{Within-COVID, models differ from each other, but no foundation model beats the from-scratch baseline under the pre-specified test}
\label{ssec:r1_within_covid}

We first asked whether foundation-model embeddings carry any patient-level
mortality signal within COVID-19. For each model, frozen embeddings of the same
raw platelet transcriptomes were classified with logistic regression under
repeated patient-level cross-validation (10 repeats of 5-fold
\texttt{StratifiedGroupKFold} with \texttt{patient\_id} as the group), and each
patient's out-of-fold cell predictions were averaged into a single score before
computing the patient-level out-of-fold ROC-AUC (patient-level AUC;
Section~\ref{ssec:methods_metrics}). Three from-scratch controls ran under the
identical protocol: a PCA-200 + logistic-regression baseline on raw counts, a
null model given only each patient's cell count and disease label, and a
cohort-only model given each patient's cohort mortality rate.

Patient-level mortality signal was learnable: five of the six methods,
including the from-scratch baseline, significantly outscored the null model
(exact DeLong paired tests on the seed-ensemble scores, $p = 0.0006$--$0.023$);
the exception was \model{UCE} ($p = 0.157$; Supplementary
Note~\ref{sn:method_vs_null}). The pre-specified primary contrast
(Section~\ref{ssec:methods_cv}) asked whether the best foundation model beats
the from-scratch baseline once winner selection is accounted for. It was not
met: the observed maximum advantage (\model{Geneformer}, $+0.064$
patient-level AUC on the seed-ensemble scores) did not exceed what the
selection procedure produces under the null (patient-level label-permutation
max-statistic test, 10{,}000 permutations, $P = 0.181$; recentered-bootstrap
cross-check, $P = 0.073$; Supplementary Note~\ref{sn:winner_test}). The result
is stable: all 25 leave-one-fatal-patient-out refits returned $P > 0.05$
(range 0.141--0.288; Supplementary Note~\ref{sn:fragility}).

The model ranking underneath this null is still informative
(Table~\ref{tab:within_covid} and Figure~\ref{fig:within_covid}). \model{Geneformer} and \model{STATE} led at
0.786 $\pm$ 0.026 and 0.778 $\pm$ 0.033 (repeat-mean $\pm$ repeat-sd) and
could not be distinguished from each other ($\Delta = -0.012$, exact DeLong
$p = 0.58$). Their nominal advantages over the baseline (\model{Geneformer}
$+0.064$, $p = 0.036$; \model{STATE} $+0.052$, $p = 0.048$) are reported as
exploratory and do not survive the winner correction above. All remaining
contrasts are likewise exploratory
(Supplementary Note~\ref{sn:se_check}):
\model{scPRINT-2} placed third (0.740 $\pm$ 0.027) and could not be separated
from the baseline ($p = 0.65$), and \model{scGPT} (0.720 $\pm$ 0.047) was
statistically indistinguishable from it ($p = 0.91$). \model{UCE} placed last
(0.661 $\pm$ 0.046), below the from-scratch baseline, though the paired
contrast was not itself significant ($-0.066$, $p = 0.081$). One separation
survives selection correction: the two leaders' advantage over \model{UCE}
($+0.118$ and $+0.130$) is confirmed by a range-statistic test on the same
label permutations as the primary analysis (max-minus-min over the five
foundation models, $P = 0.011$; Supplementary Note~\ref{sn:range_test}) and
is the paper's single confirmatory secondary result. Within the exploratory
pairwise family, six contrasts survive Benjamini--Hochberg correction, all
with \model{UCE} or \model{scGPT} on the losing side ($q \leq 0.037$;
Supplementary Figure~\ref{fig:supp_pairwise}), and the two
leaders-versus-\model{UCE} gaps are the largest. The result's content is a
statement about the losing model, addressed in
Section~\ref{ssec:discussion_property}.

A separate question, whether a model's score adds information beyond the
baseline's even without ranking patients better on its own, was examined with
nested-model comparisons: a reduced model containing cohort and the PCA-200
score against a full model adding each foundation-model score in turn. No
model's increment survived the same Benjamini--Hochberg correction used
throughout, in any of three comparator configurations
(Supplementary Note~\ref{sn:nested_comparison}). The complementarity question
therefore returns the same answer as the superiority question, and the
headline statement needs no qualifier.

Three controls bound where the learnable signal does not come from. The null
model reached only 0.554, so mortality was not readable from cell counts or
disease identity alone; per-patient cell counts also correlated only weakly
with predicted scores across models (Spearman $\rho = 0.18$--$0.29$;
Supplementary Figure~\ref{fig:supp_cellcount}). The cohort-only model reached
0.486 $\pm$ 0.038, statistically indistinguishable from chance; because a
rate estimator based on add-one smoothing inverts risk ordering for small
cohorts, this number is silent on the confound question rather than
exonerating cohort (Supplementary Note~\ref{sn:cohort_baseline}). A
QC-metrics-only model scored below chance out-of-fold (0.376 as-fit), a
textbook suppression pattern: the QC features carry outcome information only
conditional on library depth, and adding each patient's cell count recovers a
correctly oriented 0.618 (Supplementary Note~\ref{sn:qc_baseline}). The
decisive test is incremental value over all of these covariates at once:
a baseline of QC metrics, cell count, and cohort identity reached 0.627, and
the pre-specified primary for this question, whether the from-scratch
platelet representation adds information over that baseline, chosen because
it involves no model selection, was met by a nested likelihood-ratio test
against an empirical permutation null ($P < 0.0001$;
Supplementary Note~\ref{sn:qc_incremental}). The five foundation models,
reported as exploratory with both readings paired, each added significant
value by the same test ($P < 0.0001$--$0.0022$, \model{UCE} weakest) and, for
all but \model{UCE}, by paired-AUC comparison ($\Delta = +0.026$ to $+0.065$,
exact DeLong $p = 0.0001$--$0.006$; \model{UCE} $+0.016$, $p = 0.072$). The
likelihood-ratio family establishes that incremental information exists; the
AUC increments quantify it; and the resolution argument applies to that
magnitude, not to the likelihood-ratio $P$-values.

Two protocol observations qualify the headline numbers. First, random-forest
classifiers scored uniformly lower than logistic regression on the same
embeddings (e.g., \model{STATE}: 0.696 versus 0.778), so we report logistic
regression throughout and retain random forest only for calibration analysis
(Section~\ref{ssec:r5_calibration}). Second, per-seed variation was large enough
to change the apparent ranking: \model{STATE}'s patient-level AUC ranged from
0.711 to 0.815 across the ten repeats, and a single 5-fold split would have
reported any point in that range as ``the'' result. We therefore treat the
repeat-mean, not any single split, as the primary estimate throughout, with
inferential quantities computed on the seed-ensemble scores
(Section~\ref{ssec:methods_cv}).

The within-COVID picture is thus one of inter-model differences without a
detectable pretraining advantage. Signal is present and learnable by nearly
every method, and it is not reducible to library quality, cell counts, or
cohort membership: platelet representations carry outcome information beyond
all three. Yet no foundation-model representation carries information beyond
what a from-scratch PCA of the same data already provides, not in
discrimination and not in nested complementarity. The two leading models
separate from \model{UCE} under a confirmatory test, a statement about the
weakest model rather than the strongest. At 25 fatal events this benchmark
resolves advantages of roughly 0.10 patient-level AUC or larger between two
different representations (Section~\ref{ssec:discussion_property}): the
nominal 0.05--0.06 advantages sit below that resolution, and detecting effects
of that size would require on the order of $10^3$ patients. Whether any signal
transfers beyond COVID-19 is addressed next
(Sections~\ref{ssec:r2_within_sepsis}--\ref{ssec:r3_cross_disease}).

% Table: within-COVID main table (tables_v3.md Table 3; sources cited therein)
\begin{table}[t]
\centering
\caption{Within-COVID-19 patient-level mortality prediction (240 patients, 25
fatal; repeated 10$\times$5-fold cross-validation, logistic regression).
Repeat-mean $\pm$ repeat-sd is the primary reporting unit; the ensemble-OOF
AUC averages each patient's scores across the ten repeats before scoring
once. Pairwise contrasts are exploratory nominal exact-DeLong $p$-values. The
pre-specified primary contrast (best foundation model versus PCA-200) was not
met: winner-corrected permutation $P = 0.181$, recentered-bootstrap
cross-check $P = 0.073$. The single confirmatory secondary result is the two
leaders' separation from \model{UCE} (range-statistic permutation
$P = 0.011$).}
\label{tab:within_covid}
\small
\begin{tabular}{lcccc}
\toprule
\textbf{Method} & \textbf{Repeat-mean $\pm$ sd} & \textbf{Ensemble-OOF AUC} & \textbf{vs PCA-200 (DeLong $p$)} & \textbf{vs null (DeLong $p$)} \\
\midrule
\model{Geneformer} & 0.786 $\pm$ 0.026 & 0.807 & 0.036 & 0.0006 \\
\model{STATE} & 0.778 $\pm$ 0.033 & 0.795 & 0.048 & 0.0014 \\
\model{scPRINT-2} & 0.740 $\pm$ 0.027 & 0.756 & 0.655 & 0.011 \\
PCA-200+LR (baseline) & 0.725 $\pm$ 0.033 & 0.742 & --- & 0.013 \\
\model{scGPT} & 0.720 $\pm$ 0.047 & 0.739 & 0.909 & 0.023 \\
\model{UCE} & 0.661 $\pm$ 0.046 & 0.677 & 0.081 & 0.157 \\
\midrule
null (cell count + disease) & 0.554 & --- & --- & --- \\
cohort-only & 0.486 $\pm$ 0.038 & --- & --- & --- \\
QC metrics + cell count (as-fit) & 0.618 & --- & --- & --- \\
\bottomrule
\end{tabular}
\end{table}

\begin{figure}[t]
\centering
\includegraphics[width=\textwidth]{figures/v2/Figure2_within_covid.pdf}
\caption{\textbf{Within-COVID-19 mortality prediction.} \textbf{(A)}
Patient-level AUC (repeat-mean $\pm$ repeat-sd over 10$\times$5-fold
cross-validation) for the five foundation models, the PCA-200+LR baseline
(open circle), and three controls (gray squares: null = cell count +
disease; cohort-only; QC + cell count). Bracket: the pre-specified primary
contrast (best foundation model versus PCA-200), winner-corrected permutation
$P = 0.181$ (n.s.); the confirmatory leaders-versus-\model{UCE} range
statistic, $P = 0.011$. \textbf{(B)} Incremental value over a
QC + cell count + cohort baseline (AUC 0.627): points are exact-DeLong
$\Delta$AUC with 1.96$\times$SE whiskers, annotated with DeLong $p$ and the
nested likelihood-ratio permutation $P$; the PCA-200 row is the
pre-specified primary for this question ($P < 0.0001$).}
\label{fig:within_covid}
\end{figure}

% R2 — within-sepsis mortality prediction is null at late sampling windows (draft 2026-09-07)
\subsection{Within-sepsis, no model predicts mortality above chance at late sampling windows}
\label{ssec:r2_within_sepsis}

We next ran the identical protocol on the sepsis arm (91 patients, 17 fatal,
28-day mortality; Section~\ref{ssec:methods_cohort}). Sepsis is
life-threatening organ dysfunction caused by a dysregulated host response to
infection \citep{singer2016sepsis3} and accounts for roughly 11 million
deaths annually worldwide \citep{rudd2020global}. This is, to our
knowledge, the largest platelet single-cell sepsis cohort with per-patient
outcome labels assembled to date, and the first evaluated under patient-level
cross-validation. No model exceeded chance
(Table~\ref{tab:within_sepsis} and Figure~\ref{fig:within_sepsis}).
\model{UCE} scored highest (patient-level AUC 0.548 $\pm$ 0.038), followed by
\model{scPRINT-2} (0.523 $\pm$ 0.049), \model{Geneformer} (0.515 $\pm$ 0.043),
and \model{STATE} (0.511 $\pm$ 0.040), with \model{scGPT} (0.449 $\pm$ 0.052)
and the PCA-200 baseline (0.428 $\pm$ 0.082) both below chance. Under
patient-level paired bootstrap with multi-comparison control, no pairwise
difference between models survived:
the only nominal contrast was \model{UCE} over \model{scGPT} ($+0.101$,
95\% CI [0.001, 0.197], $p = 0.047$, not surviving correction), and no model
separated from the PCA-200 baseline (all 95\% confidence intervals include
zero; smallest nominal $p = 0.079$ for \model{UCE}). In other words, the
within-sepsis comparison is a complete null: not only does every method sit at
chance, no ordering among them can be defended at this sample size. At 17
fatal events the benchmark resolves only large effects: 80\% power to detect
a single predictor requires a patient-level AUC of about 0.71 or higher, so
this null bounds rather than excludes smaller signal.

The cohort's sampling times complicate any stronger reading
(Section~\ref{ssec:limitations}). Draws occurred a median of 16 days after
admission (range 3--121), and fatal cases were drawn earlier than survivors
(median 10 versus 18.5 days), so sampling time is associated with the outcome
and could, in principle, masquerade as signal. The measured result is a null
despite that confound working in the models' favor. The literature context
sharpens the interpretation: outcome-predictive platelet transcriptomic signals
in sepsis have so far been documented only in early windows, within hours of
diagnosis or within 72 hours of ICU admission
\citep{qiu2021sepsis, middleton2019sepsis, middleton2020nets}, against a
background of sepsis host-response heterogeneity established by single-cell
and transcriptomic studies \citep{reyes2020immune, davenport2016genomic,
sweeney2018unsupervised}. Our measurement is the first at
draws beyond the first days of illness, and it bounds when that signal remains
predictable: by the second week after admission, neither foundation-model
embeddings nor simple baselines carried usable patient-level mortality signal.
Because the early-window evidence comes from different cohorts and platforms,
this comparison is between studies; a within-cohort temporal test would
require serial draws we do not have.

% Table: within-sepsis (P0_onepager.md Sec.3b)
\begin{table}[t]
\centering
\caption{Within-sepsis mortality prediction (91 GSE279452 patients, 17 fatal,
28-day endpoint; same repeated-CV protocol as
Table~\ref{tab:within_covid}). No pairwise contrast survives
multiple-comparison control; the only nominal one was \model{UCE} over
\model{scGPT} ($p = 0.047$, not corrected-significant), and no model
separates from the PCA-200 baseline (all 95\% confidence intervals include
zero). At 17 fatal events the design resolves only AUC $\gtrsim 0.71$ at
80\% power, so the null bounds rather than excludes smaller signal.}
\label{tab:within_sepsis}
\small
\begin{tabular}{lc}
\toprule
\textbf{Method} & \textbf{Patient-level AUC (repeat-mean $\pm$ sd)} \\
\midrule
\model{UCE} & 0.548 $\pm$ 0.038 \\
\model{scPRINT-2} & 0.523 $\pm$ 0.049 \\
\model{Geneformer} & 0.515 $\pm$ 0.043 \\
\model{STATE} & 0.511 $\pm$ 0.040 \\
\model{scGPT} & 0.449 $\pm$ 0.052 \\
PCA-200+LR (baseline) & 0.428 $\pm$ 0.082 \\
\bottomrule
\end{tabular}
\end{table}

\begin{figure}[t]
\centering
\includegraphics[width=0.6\textwidth]{figures/v2/Figure3_within_sepsis.pdf}
\caption{\textbf{Within-sepsis mortality prediction is null at late sampling
windows.} Distribution of per-repeat patient-level AUCs (10$\times$5-fold
cross-validation; violin, per-repeat points, and mean) for the five
foundation models and the PCA-200+LR baseline on the 91-patient/17-fatal
GSE279452 arm (28-day mortality, median draw 16 days post admission). Every
method sits at or below chance, and no model ordering is defensible at this
sample size (Table~\ref{tab:within_sepsis}).}
\label{fig:within_sepsis}
\end{figure}

% R3 — cross-disease transfer and the adaptation curve (synced to P0 round-17 state, 2026-09-09)
\subsection{Cross-disease transfer is unresolvable in one direction and rank-indeterminate in the other}
\label{ssec:r3_cross_disease}

We then tested whether representations learned in one disease carry to the
other, training on COVID-19 and testing on sepsis and vice versa under the
primary severe-as-survived label definition, with patient-level scoring and
patient-clustered bootstrap intervals (Section~\ref{ssec:methods_transfer}).
Forward transfer (COVID$\rightarrow$sepsis; 91 sepsis patients, 17 fatal) is
uninformative by construction: Section~\ref{ssec:r2_within_sepsis} establishes
that mortality is not learnable in this pool even with in-domain training, so
a transfer measurement into it re-measures that null rather than transfer
ability. Point estimates spanned 0.644 (\model{Geneformer}) to 0.684
(\model{UCE}), with bootstrap confidence intervals overlapping throughout
(Table~\ref{tab:transfer} and Figure~\ref{fig:transfer}). Reverse transfer
(sepsis$\rightarrow$COVID; 96 sepsis training patients, 240 COVID-19 test
patients, 25 fatal) tests the direction where signal exists and gave higher
point estimates throughout: \model{STATE} 0.721 (95\% CI [0.621, 0.809]) and
\model{Geneformer} 0.712 (95\% CI [0.612, 0.814]) with logistic regression
(0.711 and 0.636 with random forest), followed by \model{UCE} (0.693/0.607),
\model{scGPT} (0.642/0.622), and \model{scPRINT-2} (0.619/0.517) (logistic
regression/random forest), with all bootstrap intervals overlapping. No
reverse-direction ranking can be defended either: the severe-excluded
sensitivity analysis reverses the top two models (\model{Geneformer} 0.775,
\model{STATE} 0.765), and adding only five training patients moves
\model{UCE} by $+0.086$ AUC between the 91- and 96-patient training pools
(Table~\ref{tab:transfer}). Transfer estimates are therefore
not identifiable at this scale, which is itself the strongest evidence the
experiment provides. The forward/reverse asymmetry is a property of the two
target tasks, one learnable and one not, not of transfer direction.

To ask whether transfer can be recovered with target-disease supervision, we
measured the adaptation curve: COVID-trained classifiers received increasing
numbers of sepsis training cells and were scored on a held-out 32-patient
sepsis test set, a one-third stratified holdout of the 96-patient labeled
sepsis pool (Section~\ref{ssec:methods_transfer}). Even with all
target-disease training data, all five foundation models and the PCA-200
baseline remained between 0.40 and 0.62 patient-level AUC (zero-shot values
lower still), with no foundation model separating from the baseline at any
training size (Figure~\ref{fig:adaptation_curve}). We report this as
unresolvable rather than null: at roughly 3 fatal events in 32 test patients,
the 95\% interval around any patient-level AUC spans about 0.7, so the
observed range is consistent with a single value plus noise, and no
rescue-or-not conclusion is identifiable at this test-set size. Two
comparability notes accompany this result. First, cross-disease
transfer on the 5-patient GSE167363 sepsis test set (2 fatal) produced
substantially higher point estimates (up to 1.00 patient-level AUC); given
the test-set size, we report those values
only as a labeled small-sample reference (Supplementary Note~\ref{sn:gse167363_testset}). Second, two additional
models evaluated under a different test split (TranscriptFormer and STACK) are
excluded from this comparison for the same reason and are noted in the same
supplementary record.

Together, the transfer experiments provide no evidence for a transferable
inflammatory-mortality representation on this underrepresented cell type:
forward transfer is uninformative, reverse transfer is rank-indeterminate,
and the adaptation curve is below the resolution of its test set.

% Table: cross-disease transfer (tables_v3.md Table 5)
\begin{table}[t]
\centering
\caption{Cross-disease transfer under the primary severe-as-survived label
definition (patient-level scoring; 95\% patient-clustered-bootstrap intervals,
200 iterations). Forward: train on COVID-19 (240 patients), test on the
91-patient GSE279452 sepsis pool (17 fatal) --- uninformative by construction
because the target task is not learnable in-domain
(Table~\ref{tab:within_sepsis}). Reverse: train on all 96 outcome-labeled
sepsis patients (91 GSE279452 + 5 GSE167363), test on COVID-19 (240 patients,
25 fatal). The severe-excluded column is the label-variant sensitivity
analysis; it reverses the top two reverse-direction models, so no ranking is
reported.}
\label{tab:transfer}
\small
\resizebox{\textwidth}{!}{%
\begin{tabular}{llccc}
\toprule
\textbf{Direction} & \textbf{Model} & \textbf{LR AUC [95\% CI]} & \textbf{RF AUC} & \textbf{Severe-excluded LR AUC} \\
\midrule
forward (COVID$\rightarrow$sepsis) & \model{UCE} & 0.684 [0.543, 0.805] & --- & --- \\
 & \model{STATE} & 0.653 [0.490, 0.808] & --- & --- \\
 & \model{scGPT} & 0.650 [0.495, 0.780] & --- & --- \\
 & \model{scPRINT-2} & 0.647 [0.496, 0.791] & --- & --- \\
 & \model{Geneformer} & 0.644 [0.486, 0.824] & --- & --- \\
\midrule
reverse (sepsis$\rightarrow$COVID) & \model{STATE} & 0.721 [0.621, 0.809] & 0.711 & 0.765 \\
 & \model{Geneformer} & 0.712 [0.612, 0.814] & 0.636 & 0.775 \\
 & \model{UCE} & 0.693 [0.596, 0.799] & 0.607 & 0.772 \\
 & \model{scGPT} & 0.642 [0.525, 0.765] & 0.622 & 0.678 \\
 & \model{scPRINT-2} & 0.619 [0.493, 0.743] & 0.517 & 0.630 \\
\bottomrule
\end{tabular}%
}
\end{table}

\begin{figure}[t]
\centering
\includegraphics[width=\textwidth]{figures/v2/Figure5_cross_disease.pdf}
\caption{\textbf{Cross-disease transfer and the adaptation curve.}
\textbf{(A)} Bidirectional transfer per model (logistic regression): forward
(COVID$\rightarrow$sepsis, 91-patient GSE279452 test pool) versus reverse
(sepsis$\rightarrow$COVID, 96-patient training pool, 240-patient test). All
bootstrap confidence intervals overlap (Table~\ref{tab:transfer}); forward
transfer re-measures the within-sepsis null, and the reverse direction is
rank-indeterminate (the severe-excluded sensitivity analysis reverses the
top two models). \textbf{(B)} Adaptation curve: COVID-trained classifiers
given increasing numbers of sepsis training cells and scored on a held-out
32-patient sepsis test set. The gray band marks the unresolvable region at
this test-set size (approximately three fatal events; 95\% interval width
$\approx 0.7$); all methods lie inside it throughout.}
\label{fig:transfer}
\label{fig:adaptation_curve}
\end{figure}

% R4 — in-distribution positive control on monocytes (draft 2026-09-07)
\subsection{The in-distribution positive control returns no separation either}
\label{ssec:r4_monocyte}

A reviewer could dismiss the platelet results as an artifact of evaluating too
far outside the models' intended use. We therefore reran the identical pipeline
on a cell type the models were trained on: monocytes from the same source
cohorts (Ren and Stephenson COVID-19; five Qiu sepsis patients; 148
patients\footnote{One Ren survivor patient (S-S072) is not represented in
the monocyte cell set, so the monocyte cohort contains 148 rather than 149
patients and the disease-matched COVID-only subset 143 rather than 144;
fatal counts are unaffected.}, 19 fatal; Section~\ref{ssec:methods_baselines}). Monocytes are among the most
abundant immune cells in the atlases these models were pretrained on, so this
arm isolates cell-type familiarity as the single changed variable.

The positive control returned no separation at all
(Table~\ref{tab:monocyte} and Figure~\ref{fig:indist_vs_ood}). Under
patient-level paired bootstrap, no foundation model significantly outperformed
the PCA-200 baseline on monocytes (\model{STATE} $+0.043$, 95\% CI [$-0.026$,
0.120], $p = 0.23$; \model{Geneformer} $+0.022$, $p = 0.57$; \model{scPRINT-2}
$+0.018$, $p = 0.68$; \model{UCE} $+0.012$, $p = 0.80$; \model{scGPT} $+0.005$,
$p = 0.94$). Even on a cell type these models were trained on, the from-scratch
baseline is statistically sufficient. One directional pattern is worth
noting: \model{UCE}'s monocyte-versus-platelet difference is the closest to
zero of any method in both cohort pools ($-0.005$ COVID-only, $-0.014$
all-cohort, versus drops of $0.043$--$0.119$ for the other five methods) ---
the least affected by cell type, though it beats the baseline on neither.

One inversion requires careful reading. Every method, including the baseline,
scored lower on monocytes than on platelets (baseline $-0.072$; confirmed on
the disease-matched COVID-only subset, 143 patients, 17 fatal). Because the
from-scratch baseline drops by as much as the foundation models, the inversion
is a property of the task and cohort, not of pretraining: patient-level
mortality from monocytes is simply a harder measurement in these cohorts. That
platelet transcriptomes carry more mortality signal than monocyte
transcriptomes in COVID-19 is consistent with the mechanistic literature
placing platelet-driven immunothrombosis proximal to fatal disease
\citep{manne2020platelet,barrett2021platelets}. The positive control therefore narrows the out-of-distribution
objection in an unexpected direction. Even on a cell type the models were
trained on, the from-scratch baseline is statistically sufficient, and with
the within-COVID primary contrast also returning null under the
winner-corrected test (Section~\ref{ssec:r1_within_covid}), the pipeline
resolved no pretraining advantage over the baseline in any cell of the
design. The platelet results can therefore not be attributed to evaluating
too far outside the models' intended use: the same comparison is null
in-distribution. Equally, the control cannot certify that this pipeline
resolves foundation-model advantages anywhere, and we report it with that
qualification (Section~\ref{ssec:limitations}).

% Table: monocyte positive control (P0_onepager.md Sec.3c)
\begin{table}[t]
\centering
\caption{In-distribution positive control on monocytes from the same source
cohorts (148 patients, 19 fatal; same repeated-CV protocol). Patient-level
AUC point estimates with the paired-bootstrap contrast against the PCA-200
baseline; no foundation model separates from the baseline
(\model{STATE} $+0.043$, 95\% CI [$-0.026$, 0.120]). Every method, including
the baseline, scores lower on monocytes than on platelets (baseline
$-0.072$), a task property, not a model property.}
\label{tab:monocyte}
\small
\begin{tabular}{lcc}
\toprule
\textbf{Method} & \textbf{Patient-level AUC} & \textbf{vs PCA-200 ($\Delta$, $p$)} \\
\midrule
\model{STATE} & 0.696 & $+0.043$, $p = 0.23$ \\
\model{scGPT} & 0.677 & $+0.005$, $p = 0.94$ \\
\model{Geneformer} & 0.674 & $+0.022$, $p = 0.57$ \\
\model{scPRINT-2} & 0.672 & $+0.018$, $p = 0.68$ \\
PCA-200+LR (baseline) & 0.654 & --- \\
\model{UCE} & 0.648 & $+0.012$, $p = 0.80$ \\
\bottomrule
\end{tabular}
\end{table}

\begin{figure}[t]
\centering
\includegraphics[width=\textwidth]{figures/v2/Figure4_indist_vs_ood.pdf}
\caption{\textbf{In-distribution (monocytes) versus underrepresented
(platelets) mortality prediction.} \textbf{(A)} Patient-level AUC per method
on the 148-patient/19-fatal monocyte control (all cohorts, including five Qiu
sepsis patients) against the within-COVID-19 platelet arm: no foundation
model separates from the PCA-200 baseline on either cell type, and every
method scores lower on the in-distribution cell type. \textbf{(B)}--\textbf{(C)}
The COVID-only matched subset (143 patients, 17 fatal) confirms the inversion
is not a cohort-composition artifact.}
\label{fig:indist_vs_ood}
\end{figure}

% R5 — calibration and robustness axes (draft 2026-09-07)
\subsection{Calibration favors STATE, and random forest probabilities are better calibrated than logistic-regression scores}
\label{ssec:r5_calibration}

Discrimination is not the only property that decides whether a risk estimate
can be acted on: the emitted probability must also mean what it says. We
computed patient-level calibration on the within-COVID out-of-fold predictions
(Figure~\ref{fig:calibration}). Among logistic-regression classifiers,
\model{STATE} was the best calibrated (Brier 0.118, expected calibration error
0.168), ahead of \model{Geneformer} (0.119, 0.190) and \model{scPRINT-2}
(0.121, 0.181), with \model{UCE} (0.162, 0.243) and \model{scGPT} (0.165,
0.257) the least calibrated. Random-forest probabilities were better
calibrated than logistic-regression scores for every model (ECE 0.054--0.080
versus 0.168--0.257), reflecting the conservatism of tree-averaged
probabilities \citep{niculescu2005predicting, naeini2015obtaining}, so where calibrated risk estimates matter for deployment we
recommend random forest on the leading embedding (\model{STATE}-RF: Brier
0.095, ECE 0.079) over any logistic-regression variant.

Two further secondary axes complete the multi-axis profile. First,
\model{UCE} was the most resistant to embedding-dimension dropout: its AUC did
not degrade under 70\% dimension zeroing, indicating a highly distributed
representation (Supplementary Figure~\ref{fig:supp_dropout}; computed on the
legacy analysis pool, an earlier cohort assembly predating the canonical
240-patient cohort, Section~\ref{ssec:limitations}). Second,
leave-one-cohort-out evaluation on the analysis cohort measured
generalization to cohorts that contain fatal cases, an outcome-selected
estimand: folds are computable only when the held-out cohort contains both
classes, so survivor-only cohorts contribute to training throughout and are
never scored out-of-fold (Supplementary Figure~\ref{fig:supp_lobo}).
Because the five computable folds range from 13 to 90 patients and from 1 to
9 fatal events, we report the AUC over concatenated held-out predictions
(206 patients, 25 fatal events) rather than a mean of per-fold AUCs: the top
three models sit within 0.005 AUC of each other (\model{STATE} 0.753,
\model{scPRINT-2} 0.751, \model{Geneformer} 0.748), ahead of \model{UCE}
(0.681) and \model{scGPT} (0.638). Per-fold AUCs and the per-cohort counts
are given in Supplementary Table~\ref{tab:lobo_folds}. By fold-to-fold spread,
\model{scGPT} was the most cohort-consistent ($\sigma = 0.123$) and
\model{UCE} the least ($\sigma = 0.152$). No single model leads on every
axis; the choice among them is deployment-specific
(Section~\ref{ssec:discussion_guidance}).

\begin{figure}[t]
\centering
\includegraphics[width=\textwidth]{figures/v2/Figure7_calibration.pdf}
\caption{\textbf{Calibration of patient-level mortality predictions within
COVID-19.} Reliability diagrams (10 equal-width bins; point area proportional
to bin population) for \textbf{(A)} logistic-regression and \textbf{(B)}
random-forest classifiers on each model's out-of-fold predictions.
\model{STATE} is the best-calibrated logistic-regression model (Brier 0.118,
ECE 0.168); random-forest probabilities are better calibrated for every
model (ECE 0.054--0.080), with \model{STATE}-RF at Brier 0.095, ECE 0.079.}
\label{fig:calibration}
\end{figure}

% R6 — unified masking-based gene attribution + cross-model GSEA comparison
% (masking thread completed 2026-09-11: 5 models, full-cohort runs, GSEA,
% joint-occlusion dose control; all numbers verified against
% scripts/31_donor_level_revision/results_gene_attribution_gsea/)
\subsection{The five models prioritize different biological programs for the same prediction, and pathway-level signal is decoupled from predictive performance}
\label{ssec:r6_attribution_gsea}

The analyses above ask how well each model predicts. Here we ask a different
question: when the models' representations are mapped back to genes, do they
use the \emph{same biology} to predict mortality? We built a unified
masking-based (leave-one-gene-out) attribution pipeline
(Section~\ref{ssec:methods_masking}; Figure~\ref{fig:masking_design}): for
every cell, the same top-256
expressed candidate genes were selected from the intersection of the models'
vocabularies; each gene was occluded in turn through the model's native
attention mechanism, and the change in the mortality-classifier score,
$\Delta(g, \text{cell}) = s_\text{full} - s_\text{masked}(g)$, was averaged
across all 10{,}737 cells into a per-gene mortality attribution for each
model. The classifier is a single in-sample logistic regression per model
($C = 0.3$) that defines each model's mortality direction; its AUC is an
internal direction-quality diagnostic, not an out-of-fold performance
estimate, and is not comparable to the cross-validated AUCs of
Section~\ref{ssec:r1_within_covid}.

\begin{figure}[t]
\centering
\includegraphics[width=\textwidth]{figures/v2/Figure_masking_design.png}
\caption{\textbf{Unified masking-based gene attribution across the five
foundation models.} Every model receives the same cells and, per cell, the
same top-256 expressed candidate genes from the intersection of the five
model vocabularies, and an intact-cell mortality-classifier score is
established before any gene is removed (Stages 1--2). Each candidate gene is
then occluded in turn and the change in score,
$\Delta(g, \text{cell}) = s_\text{full} - s_\text{masked}(g)$, is recorded
(Stage 3). Occlusion uses each architecture's native attention mechanism ---
attention-mask flag for \model{Geneformer}, physical sequence deletion for
\model{STATE} (the only model whose mask flag does not reach attention),
padding-mask flag for \model{scGPT} and \model{UCE}, and an
inverted-polarity mask flag for \model{scPRINT-2} --- so the measured
quantity is identical across architectures (Stage 4). Per-gene deltas are
averaged across the 10{,}737 cells into a ranked list per model (Stage 5)
and taken to prerank GSEA for the cross-model pathway comparison (Stage 6).}
\label{fig:masking_design}
\end{figure}

\textbf{The five models agree on the task but not on the genes.} Each
model's classifier recovered the mortality direction from its own
256-gene-restricted embeddings (patient-level in-sample AUC 0.760--0.990;
Table~\ref{tab:attribution_diagnostics}), yet the per-gene attributions were
only weakly concordant across models: pairwise Spearman rank correlations of
the full ranked lists ranged from $-0.03$ to $0.40$ on the 3{,}917 genes
ranked by all five models. \model{Geneformer}'s top genes were dominated by
ribosomal proteins (RPLP2, RPL24, RPL18A) with TIMP1; \model{scGPT}'s by
glycolytic and metabolic enzymes (ALDOA, G6PD, PKM, SHMT2); \model{UCE}'s by
myeloid and tissue-remodeling genes (AXL, C1QA, PLCG2, IL1R2);
\model{STATE}'s by TIMP1, RSU1, TNFSF13B and PFN1; \model{scPRINT-2}'s by
housekeeping genes with no coherent theme. Full 5-way sign agreement is more
common than the pairwise correlations alone suggest: 531 of the 3{,}917
commonly-ranked genes (365 positive, 166 negative) carry the same sign in
all five models, exceeding the $\sim$272 expected under the models'
independent marginal sign rates. AMFR, SHMT2 and ARID5B are among the
positively-agreeing genes (AMFR and ARID5B ranking 9th--17th in
\model{STATE} and \model{UCE}, SHMT2 5th in \model{scGPT} and 14th in
\model{STATE}), and CCL5 is negatively attributed in all five (ODC1 in
four, the fifth indistinguishable from zero) --- but several prominent
platelet genes carry opposite signs across models: PPBP is among the most
negative genes for \model{Geneformer} and \model{scGPT} yet mildly positive
for \model{STATE} and \model{UCE}, and TMSB4X is positive for
\model{Geneformer} and \model{scGPT} but the fifth-most negative gene for
\model{STATE} (Figure~\ref{fig:attribution_gsea}A;
Supplementary Note~\ref{sn:attribution_gsea}).

\textbf{Pathway-level signal is model-specific, not shared.} Prerank GSEA
(Hallmark, KEGG, Reactome; FDR $q < 0.25$) on each model's ranked list found
\emph{no pathway significant in four or more of the five models}. The shared
structure that does exist is pairwise or three-way, and frequently
sign-discordant (Figure~\ref{fig:attribution_gsea}B): oxidative
phosphorylation is the one program enriched positively in three models
(\model{Geneformer}, \model{scGPT}, \model{UCE}); interferon-alpha and
interferon-gamma responses are positively enriched in \model{Geneformer}
\emph{and} \model{scGPT} but negatively in \model{UCE}; the
TNF-$\alpha$/NF-$\kappa$B and allograft-rejection programs are
significantly \emph{negative} in \model{scGPT} and \model{UCE} and trend
the same (negative) direction, non-significantly, in \model{Geneformer}
($NES=-1.46$ and $-1.22$); and immunoregulatory signaling is the true
sign-flip among the three --- significantly \emph{positive} in
\model{Geneformer} (Reactome, Immunoregulatory Interactions Between A
Lymphoid And A Non-Lymphoid Cell, $NES=2.09$) yet significantly negative in
\model{scGPT} and \model{UCE}. \model{UCE} additionally
carries its own axis --- extracellular-matrix organization and
complement/coagulation --- not shared with any other model. \model{scGPT}
and \model{UCE} produced broad pathway profiles (28--63 and 23--70
significant terms respectively in Hallmark/KEGG, and 257 and 82 in
Reactome), \model{Geneformer} and \model{STATE} narrow ones (2--23 and
2--7), and \model{scPRINT-2} none at all (0 of 48 Hallmark, 0 of 252 KEGG,
0 of 864 Reactome terms).

\textbf{Predictive performance and pathway signal are decoupled.}
\model{scPRINT-2}'s mortality classifier is the second-best of the five
(patient-level in-sample AUC 0.963) yet has zero significant pathways, while
\model{UCE}'s is the weakest (0.760) yet has one of the broadest pathway
signatures (Table~\ref{tab:attribution_diagnostics}). We adjudicated
\model{scPRINT-2}'s null directly. Its single-gene attributions have a
noise-to-signal ratio of 9.4 (median within-gene SD over across-gene SD,
versus 0.4--1.6 for the other models), raising the possibility of redundant
or compensatory gene encoding masking a real program. A joint-occlusion
dose control rules this out: jointly masking the top-10, top-30, or top-128
attributed genes (up to half of each cell's candidate sentence) moved the
score by no more than masking matched random gene sets (paired tests,
$n = 10{,}736$ cells, $P = 0.21$--$0.93$ across tiers and polarities), while
the naive additive prediction from the single-gene attributions was 55--750
times larger. The classifier's mortality direction in \model{scPRINT-2} is
therefore decoupled from gene identity at every occlusion granularity
tested; whether this reflects genuinely diffuse biology or a limitation of
gene-token occlusion for this architecture cannot be resolved with the
present method and would require an independent attribution method as a
cross-check.

\textbf{Caveats.} Three scope notes apply to this section. First, the
candidate gene universe differed across models by construction: it was the
vocabulary intersection of the models built up to each run (11{,}984 genes
for \model{Geneformer} and \model{STATE}, 11{,}977 after intersecting
\model{scGPT}'s vocabulary, and progressively smaller once \model{UCE}'s and
\model{scPRINT-2}'s vocabularies were intersected); all cross-model
comparisons here are therefore computed post hoc on the 3{,}917 genes ranked
by all five models, and pathway coverage differences partly reflect these
universe differences. Second, \model{UCE}'s released inference path passes a
float (not boolean) padding mask to its transformer, which PyTorch treats as
an additive attention bias rather than an exclusion; occlusion is only
defined for \model{UCE} after a one-line boolean cast, which we applied
uniformly to full and occluded inputs --- its attributions thus describe a
minimally corrected variant of the released model. Third, cross-model
comparisons are rank-based only: each classifier's logit scale is its own,
so attribution magnitudes are not comparable across models.

% ---- Table (inline; numbers verified 2026-09-11 against
% results_gene_attribution_gsea/<model>_masking/classifier_diagnostics.csv,
% *_ranked.csv, and the three *_gsea_*.csv per model) ----
\begin{table}[htbp]
\centering
\caption{Masking-based gene attribution: per-model classifier diagnostics
(in-sample, 256-gene-restricted embeddings, canonical within-COVID cohort of
240 patients / 25 fatal, 10{,}737 cells) and GSEA breadth (FDR $q < 0.25$).
Noise/signal = median within-gene attribution SD divided by the across-gene
SD of mean attributions. AUCs are internal direction-quality diagnostics,
not out-of-fold performance estimates.}
\label{tab:attribution_diagnostics}
\small
\begin{tabular}{lcccccc}
\hline
Model & Cell AUC & Patient AUC & Ranked genes & Noise/signal & Sig. terms (H/K/R) \\
\hline
\model{Geneformer} & 0.893 & 0.990 & 5{,}038 & 0.78 & 6 / 23 / 2 \\
\model{STATE} & 0.844 & 0.933 & 5{,}000 & 0.39 & 5 / 7 / 2 \\
\model{scGPT} & 0.802 & 0.933 & 4{,}987 & 1.62 & 28 / 63 / 257 \\
\model{UCE} & 0.680 & 0.760 & 4{,}850 & 0.44 & 23 / 70 / 82 \\
\model{scPRINT-2} & 0.857 & 0.963 & 4{,}850 & 9.43 & 0 / 0 / 0 \\
\hline
\end{tabular}
\end{table}

% ---- Figure block (PDF at reports/manuscript/figures/v2/Figure_gene_attribution_gsea.pdf;
% numbering provisional, assigned at P2 assembly) ----
\begin{figure}[htbp]
\centering
\includegraphics[width=\textwidth]{figures/v2/Figure_gene_attribution_gsea.pdf}
\caption{\textbf{Masking-based gene attribution and pathway emphasis across
five foundation models.} \textbf{(A)} Union of each model's top-7
mortality-attributed genes; color is the gene's rank within that model's own
ranked list (darker = more important), and $+$/$-$ marks the attribution
sign. \textbf{(B)} GSEA NES for pathways significant (FDR $q < 0.25$) in at
least two models, plus the strongest model-specific term of each model that
has one (below the divider); dots mark significance. No pathway is significant in four or more
models; \model{scPRINT-2}'s empty column is its complete GSEA null,
adjudicated by the joint-occlusion dose control as occlusion insensitivity
of the architecture (redundant encoding excluded). Comparisons are
rank-based only.}
\label{fig:attribution_gsea}
\end{figure}

% ============================================================================
% DISCUSSION
% ============================================================================
\section{Discussion}
\label{sec:discussion}

% Discussion D1 — which property of a foundation model predicts patient-state representation? (synced to P0 round-17 state, 2026-09-09)
\subsection{Which model property tracks patient-state representation?}
\label{ssec:discussion_property}

Architecture family is not identifiable in these data: although the
within-family spread among bidirectional transformer encoders (from
\model{Geneformer} 0.786 to \model{UCE} 0.661) exceeds the gap to the only
generative model, every between-family contrast is small against the roughly
0.10-AUC resolution of a 25-event benchmark, and the strongest between-family
gaps (\model{STATE}$-$\model{scGPT} $+0.055$;
\model{Geneformer}$-$\model{scGPT} $+0.068$; exact DeLong $p = 0.014$ and
$0.008$, in agreement with the paired bootstrap;
Supplementary Note~\ref{sn:se_check}) sit exactly at that resolution's
edge. We therefore do not claim architecture-level attribution.

The one confirmatory separation, the two leaders versus \model{UCE}
(range-statistic $P = 0.011$; Supplementary Note~\ref{sn:range_test}), is a
statement about the losing model before it is a statement about the winning
ones. \model{UCE} is the only method that never significantly outscored the
count-plus-disease null model ($p = 0.157$;
Supplementary Note~\ref{sn:method_vs_null}), so the content of the gap is
that \model{UCE} is unsuited to this task, not that the leaders carry a
pretraining advantage, a claim the primary test rejects
(Section~\ref{ssec:r1_within_covid}). The natural attribution hypothesis for
such a gap, blood and immune representation in the pretraining corpus, was
measured directly and is falsified on both ends
(Supplementary Note~\ref{sn:corpus_platelet}). The three models whose
platelet content can be numerically bounded against public corpus records
(\model{scPRINT-2} 0.120--0.125\%, \model{scGPT} 0.125\%, \model{UCE}
0.120\%) share essentially identical fractions yet span 0.079 patient-level
AUC, 63\% of the entire between-model spread. Of the two leaders,
\model{STATE}'s corpus status is itself bounded and is the lowest documented
of all --- about two-thirds of its corpus (the scBaseCount and Tahoe
components) is platelet-free at both the annotation and quality-control
layers, and the remainder shares the Census bound, capping its platelet
content below 0.05\% (Table~\ref{tab:models}) --- while \model{Geneformer}
is the one model with no cell-type-level corpus composition in any published
material. A leader at a documented near-zero and a laggard at 0.120\%
produce the same ordering as one leader with a documented corpus and one
without: wherever corpus platelet content can actually be checked, it does
not track performance.

Three further facts bound any property claim. First, corpus metrics must be
read in one unit, and none exists across all five models: the fraction of
pretraining cells from blood and bone marrow is documented as dataset counts
for \model{Geneformer} (49 of 561 source datasets, under a 25\%-per-tissue
cap), as a sample fraction for \model{UCE} (6.6\%), and not at all for the
rest (Supplementary Table~\ref{tab:corpus_fractions}). Second,
\model{scPRINT-2} contradicts any simple identity-objective story: its
objective includes explicit cell-label prediction, the attribute assigned to
the losing model, yet it ranks third, above the baseline on point estimates;
and the from-scratch PCA-200 baseline, with no pretraining at all, lands
mid-table and is statistically inseparable from two of the five models.
Third, the ordering does not persist on the in-distribution monocyte control:
\model{scGPT} moves from fourth on platelets to second on monocytes
(Section~\ref{ssec:r4_monocyte}), so no corpus-level constant can explain the
ranking on both cell types.

We therefore state the property question as a set of benchmark-testable
hypotheses, one of which now carries a recorded falsification. Supplementary
Table~\ref{tab:property_multiplicity} enumerates the candidate attributes we
checked (corpus blood/marrow fraction, corpus disease-condition fraction,
objective loss target, tokenization, corpus size, parameter count, embedding
dimension, vocabulary size, species coverage, and release year) and reports
how many separate the observed ranking equally well; the corpus-composition
hypothesis above remains in the set as an untested candidate, now known to be
falsified where it is measurable and untestable where it would matter.
Resolving any such property at the 0.03--0.05 AUC scale where attribution
lives requires on the order of $10^3$ patients at this event rate; the
present benchmark is designed so that any new model or new candidate property
can be inserted and tested under the identical protocol.

\subsection{Limitations}
\label{ssec:limitations}

Every null in this benchmark is an upper bound set by a finite design, not a
proof of absence, and the bounds are reported as such: at 25 fatal events the
within-COVID arm resolves advantages of roughly 0.10 patient-level AUC between
two representations, and the 17-event sepsis arm resolves only AUC
$\gtrsim 0.71$, so effects below those floors remain compatible with the
data. Four further limitations qualify the results. First, the two sepsis
sources use different endpoints (in-hospital death for GSE167363, 28-day
death for GSE279452), and within GSE279452 the fatal cases were drawn earlier
than survivors (median 10 versus 18.5 days), so sampling time is associated
with the outcome; the null was measured despite that association favoring the
models. Second, \model{Geneformer} publishes no cell-type-level corpus
composition at all, so corpus-composition attribution cannot be assessed for
it; and across all five models, overlap between the pretraining corpora and
this evaluation cohort cannot be excluded without corpus study-level
manifests --- a structural problem this field has already flagged for
benchmark design \citep{kedzierska2025zeroshot}.
Third, the LOBO axis is an outcome-selected estimand: folds are computable
only for cohorts containing at least one fatal and one survivor patient, so
all-survivor cohorts (34 patients) contribute to training throughout and are
never scored out-of-fold; LOBO therefore measures generalization to
fatal-case-containing cohorts only. Fourth, most source studies do not report
per-patient sex, age, or ancestry, so demographic bias cannot be audited in
these data --- a field-level gap \citep{corpas2025diversity, pot2019gender},
not a property of this benchmark.

The interpretability results carry their own caveats
(Section~\ref{ssec:r6_attribution_gsea}): the attribution classifiers are
in-sample direction definitions, not out-of-fold performance estimates; the
candidate gene universes differed across models by construction, so
pathway-coverage differences partly reflect universe differences;
\model{UCE}'s attributions describe a minimally corrected variant of the
released model; and occlusion is one attribution family --- for
\model{scPRINT-2}, whether the empty pathway profile reflects diffuse biology
or a limitation of gene-token occlusion remains open pending an independent
attribution method. One secondary robustness result also predates the
canonical cohort: the embedding-dimension dropout analysis
(Section~\ref{ssec:r5_calibration}) was computed on the legacy analysis pool,
an earlier cohort assembly that predates the canonical 240-patient cohort,
and is reported as a secondary observation rather than a headline axis.

\subsection{Deployment guidance}
\label{ssec:discussion_guidance}

For patient-level outcome prediction on this cell type at this event count,
the defensible default is the from-scratch PCA + logistic-regression
baseline: it is statistically indistinguishable from every foundation model
in every design cell tested, including the in-distribution monocyte control,
and it costs orders of magnitude less to run and audit. If a foundation model
is used regardless, the choice should be made on the secondary axes rather
than on discrimination: \model{STATE} with a random-forest head is the
best-calibrated option (Brier 0.095, ECE 0.079), while \model{scGPT}'s low
cohort-to-cohort variance should not be mistaken for robustness --- it is
low variance around the weakest pooled-LOBO mean (0.638). For future models
and future claims, the protocol itself is the contribution we offer:
pre-specify a single primary contrast, correct for winner selection and
multiplicity within each family, report the resolvable effect size of the
design alongside every AUC, and include a from-scratch baseline trained under
the identical protocol \citep{collins2024tripodai}. At this event rate, detecting the 0.03--0.05 AUC
effects that property attributions live at requires on the order of $10^3$
patients; claims based on tens of patients should be read against that scale.

\section{Conclusion}
\label{sec:conclusion}

Across a patient-level benchmark pairing mortality with an underrepresented
cell type, five published single-cell foundation models did not outperform a
from-scratch baseline under a pre-specified, selection-corrected test, and
the same null held within sepsis, across diseases, and on an in-distribution
positive control. The outcome signal itself is real and not reducible to cell
counts, library quality, or cohort membership, and the models are not
interchangeable --- they prioritize measurably different genes and pathways
for the same prediction --- but the one separation that survives correction
marks the weakest model, not a winner. At realistic clinical event counts,
pretraining bought no detectable advantage; what the field can use is the
measured resolution of such benchmarks and a protocol that makes the next
claim cheap to test.


% ============================================================================
% DATA AND CODE AVAILABILITY
% ============================================================================
\section{Data and Code Availability}

All analysis code, SLURM submitters, and the LaTeX manuscript source are
publicly available at \url{https://github.com/xqiu625/302006_updategit} under
the MIT license. The benchmark pipeline is distributed as the pip-installable
Python package \texttt{scfm-benchmark} (CLI entry point \texttt{scfm-bench}),
with documentation hosted at
\url{https://xqiu625.github.io/302006_updategit/}. Conda environment manifests
are provided under \texttt{environments/} in the repository, and a Dockerfile
is included for the core analysis environment. The platelet scRNA-seq dataset
was integrated from publicly available source studies (per-cohort patient,
event, and cell counts in Table~\ref{tab:cohort}); per-study accession
identifiers (GEO/EGA) and primary citations are listed in
\texttt{docs/cohort\_sources\_reference.md} in the repository. Pre-computed
unified-input foundation-model embeddings, the patient-level result tables,
and the masking-attribution and GSEA outputs will be archived to Zenodo with
a permanent DOI upon acceptance.

% ============================================================================
% DECLARATIONS (OUP / Briefings in Bioinformatics required sections)
% ============================================================================
\section*{Acknowledgments}

We thank the UC Riverside High-Performance Computing Center for computational resources.

\section*{Author Contributions}

\textbf{Xinru Qiu}: Conceptualization, Methodology, Software, Validation, Formal analysis,
Investigation, Data curation, Visualization, Writing --- original draft,
Writing --- review \& editing.
\textbf{Adam Godzik}: Conceptualization, Methodology, Resources, Supervision,
Project administration, Writing --- review \& editing.
Both authors read and approved the final manuscript.

\section*{Conflict of Interest}

The authors declare that they have no competing interests.

\section*{Funding}

This work received no specific grant from any funding agency in the public,
commercial, or not-for-profit sectors.

\section*{Ethics Approval}

This study is a secondary analysis of de-identified, publicly available
single-cell RNA sequencing datasets obtained from GEO and EGA under the terms of
their original consent and data-use agreements (per-study accessions listed in
\texttt{docs/cohort\_sources\_reference.md}). No new human or animal data were
collected, and no institutional review board approval was required.

\section*{Author Biographies}

\noindent\textbf{Xinru Qiu} is a computational biologist in the Division of
Biomedical Sciences at the University of California, Riverside. Her research
applies machine learning and single-cell transcriptomics to immune dysregulation
in infectious and inflammatory disease, with a focus on rigorous, leakage-free
evaluation of predictive models on clinical endpoints.

\noindent\textbf{Adam Godzik} is a Professor in the Division of Biomedical
Sciences at the University of California, Riverside. His group works on
computational structural biology, protein function prediction, and the
integration of large-scale biological data for translational applications.

% ============================================================================
% REFERENCES
% ============================================================================
\begin{thebibliography}{99}

\bibitem[Cui et al.(2024)]{cui2024scgpt}
Cui, H., Wang, C., Maan, H., et al. (2024).
\newblock scGPT: toward building a foundation model for single-cell multi-omics using generative AI.
\newblock Nature Methods, 21(8), 1470--1480. doi:10.1038/s41592-024-02201-0.

\bibitem[Theodoris et al.(2023)]{theodoris2023geneformer}
Theodoris, C. V., Xiao, L., Chopra, A., et al. (2023).
\newblock Transfer learning enables predictions in network biology.
\newblock Nature, 618(7965), 616--624. doi:10.1038/s41586-023-06139-9.

\bibitem[Bommasani et al.(2021)]{bommasani2021foundation}
Bommasani, R., Hudson, D. A., Adeli, E., et al. (2021).
\newblock On the opportunities and risks of foundation models.
\newblock arXiv preprint, 2108.07258. doi:10.48550/arXiv.2108.07258.

\bibitem[Moor et al.(2023)]{moor2023generalist}
Moor, M., Banerjee, O., Abad, Z. S. H., et al. (2023).
\newblock Foundation models for generalist medical artificial intelligence.
\newblock Nature, 616, 259--265. doi:10.1038/s41586-023-05881-4.

\bibitem[Regev et al.(2017)]{regev2017hca}
Regev, A., Teichmann, S. A., Lander, E. S., et al. (2017).
\newblock The Human Cell Atlas.
\newblock eLife, 6, e27041. doi:10.7554/eLife.27041.

\bibitem[Tabula Sapiens Consortium(2022)]{tabula2022sapiens}
The Tabula Sapiens Consortium, Jones, R. C., Karkanias, J., Krasnow, M. A., et al. (2022).
\newblock The Tabula Sapiens: a multiple-organ, single-cell transcriptomic atlas of humans.
\newblock Science, 376(6594), eabl4896. doi:10.1126/science.abl4896.

\bibitem[CZI Cell Science Program et al.(2025)]{czi2025cellxgene}
CZI Cell Science Program, Abdulla, S., Aevermann, B., et al. (2025).
\newblock CZ CELLxGENE Discover: a single-cell data platform for scalable exploration, analysis and modeling of aggregated data.
\newblock Nucleic Acids Research, 53(D1), D886--D900. doi:10.1093/nar/gkae1142.

\bibitem[Yang et al.(2022)]{yang2022scbert}
Yang, F., Wang, W., Wang, F., et al. (2022).
\newblock scBERT as a large-scale pretrained deep language model for cell type annotation of single-cell RNA-seq data.
\newblock Nature Machine Intelligence, 4, 852--866. doi:10.1038/s42256-022-00534-z.

\bibitem[Hao et al.(2024)]{hao2024scfoundation}
Hao, M., Gong, J., Zeng, X., et al. (2024).
\newblock Large-scale foundation model on single-cell transcriptomics.
\newblock Nature Methods, 21, 1481--1491. doi:10.1038/s41592-024-02305-7.

\bibitem[Wen et al.(2024)]{wen2024cellplm}
Wen, H., Tang, W., Dai, X., et al. (2024).
\newblock CellPLM: pre-training of cell language model beyond single cells.
\newblock In The Twelfth International Conference on Learning Representations (ICLR).

\bibitem[Rosen et al.(2026)]{rosen2026universal}
Rosen, Y., Roohani, Y., Agrawal, A., et al. (2026).
\newblock Universal cell embedding provides a foundation model for cell biology.
\newblock Nature, advance online publication. doi:10.1038/s41586-026-10689-z.

\bibitem[Chen \& Zou(2025)]{chen2025genept}
Chen, Y., \& Zou, J. (2025).
\newblock Simple and effective embedding model for single-cell biology built from ChatGPT.
\newblock Nature Biomedical Engineering, 9(4), 483--493. doi:10.1038/s41551-024-01284-6.

\bibitem[Tejada-Lapuerta et al.(2025)]{tejada2025nicheformer}
Tejada-Lapuerta, A., Schaar, A. C., Gutgesell, R., et al. (2025).
\newblock Nicheformer: a foundation model for single-cell and spatial omics.
\newblock Nature Methods, 22(12), 2525--2538. doi:10.1038/s41592-025-02814-z.

\bibitem[Sza{\l}ata et al.(2024)]{szalata2024transformers}
Sza{\l}ata, A., Hrovatin, K., Becker, S., et al. (2024).
\newblock Transformers in single-cell omics: a review and new perspectives.
\newblock Nature Methods, 21, 1430--1443. doi:10.1038/s41592-024-02353-z.

\bibitem[Bunne et al.(2024)]{bunne2024virtualcell}
Bunne, C., Roohani, Y., Rosen, Y., et al. (2024).
\newblock How to build the virtual cell with artificial intelligence: priorities and opportunities.
\newblock Cell, 187(25), 7045--7063. doi:10.1016/j.cell.2024.11.015.

\bibitem[Kedzierska et al.(2025)]{kedzierska2025zeroshot}
Kedzierska, K. Z., Crawford, L., Amini, A. P., \& Lu, A. X. (2025).
\newblock Zero-shot evaluation reveals limitations of single-cell foundation models.
\newblock Genome Biology, 26(1), 101. doi:10.1186/s13059-025-03574-x.

\bibitem[Liu et al.(2026)]{liu2026evaluating}
Liu, T., Li, K., Wang, Y., Li, H., \& Zhao, H. (2026).
\newblock Evaluating the utilities of foundation models in single-cell data analysis.
\newblock Advanced Science, 13(27), e14490. doi:10.1002/advs.202514490.

\bibitem[Boiarsky et al.(2023)]{boiarsky2023deepdive}
Boiarsky, R., Singh, N. M., Buendia, A., Getz, G., \& Sontag, D. (2023).
\newblock A deep dive into single-cell RNA sequencing foundation models.
\newblock bioRxiv, 2023.10.19.563100. doi:10.1101/2023.10.19.563100.

\bibitem[Ahlmann-Eltze et al.(2025)]{ahlmanneltze2025linear}
Ahlmann-Eltze, C., Huber, W., \& Anders, S. (2025).
\newblock Deep-learning-based gene perturbation effect prediction does not yet outperform simple linear baselines.
\newblock Nature Methods, 22, 1657--1661. doi:10.1038/s41592-025-02772-6.

\bibitem[Wenteler et al.(2025)]{wenteler2025perteval}
Wenteler, A., Occhetta, M., Branson, N., et al. (2025).
\newblock PertEval-scFM: benchmarking single-cell foundation models for perturbation effect prediction.
\newblock In Proceedings of the 42nd International Conference on Machine Learning, PMLR 267, 66633--66677.

\bibitem[Csendes et al.(2025)]{csendes2025benchmark}
Csendes, G., Sanz, G., Szalay, K. Z., \& Szalai, B. (2025).
\newblock Benchmarking foundation cell models for post-perturbation RNA-seq prediction.
\newblock BMC Genomics, 26, 393. doi:10.1186/s12864-025-11600-2.

\bibitem[Kapoor \& Narayanan(2023)]{kapoor2023leakage}
Kapoor, S., \& Narayanan, A. (2023).
\newblock Leakage and the reproducibility crisis in machine-learning-based science.
\newblock Patterns, 4(9), 100804. doi:10.1016/j.patter.2023.100804.

\bibitem[Whalen et al.(2022)]{whalen2022pitfalls}
Whalen, S., Schreiber, J., Noble, W. S., \& Pollard, K. S. (2022).
\newblock Navigating the pitfalls of applying machine learning in genomics.
\newblock Nature Reviews Genetics, 23(3), 169--181. doi:10.1038/s41576-021-00434-9.

\bibitem[Varma \& Simon(2006)]{varma2006bias}
Varma, S., \& Simon, R. (2006).
\newblock Bias in error estimation when using cross-validation for model selection.
\newblock BMC Bioinformatics, 7(1), 91. doi:10.1186/1471-2105-7-91.

\bibitem[Fahsbender et al.(2025)]{fahsbender2025czi}
Fahsbender, E., Andersson, A., Ash, J., et al. (2025).
\newblock Benchmarking and evaluation of AI models in biology: outcomes and recommendations from the CZI Virtual Cells Workshop.
\newblock arXiv preprint, 2507.10502v2.

\bibitem[Koh et al.(2021)]{koh2021wilds}
Koh, P. W., Sagawa, S., Marklund, H., et al. (2021).
\newblock WILDS: a benchmark of in-the-wild distribution shifts.
\newblock In Proceedings of the 38th International Conference on Machine Learning, PMLR 139, 5637--5664.

\bibitem[Rowley et al.(2011)]{rowley2011genomewide}
Rowley, J. W., Oler, A. J., Tolley, N. D., et al. (2011).
\newblock Genome-wide RNA-seq analysis of human and mouse platelet transcriptomes.
\newblock Blood, 118(14), e101--e111. doi:10.1182/blood-2011-03-339705.

\bibitem[Bray et al.(2013)]{bray2013complex}
Bray, P. F., McKenzie, S. E., Edelstein, L. C., et al. (2013).
\newblock The complex transcriptional landscape of the anucleate human platelet.
\newblock BMC Genomics, 14, 1. doi:10.1186/1471-2164-14-1.

\bibitem[Middleton et al.(2019)]{middleton2019sepsis}
Middleton, E. A., Rowley, J. W., Campbell, R. A., et al. (2019).
\newblock Sepsis alters the transcriptional and translational landscape of human and murine platelets.
\newblock Blood, 134(12), 911--923. doi:10.1182/blood.2019000067.

\bibitem[Manne et al.(2020)]{manne2020platelet}
Manne, B. K., Denorme, F., Middleton, E. A., et al. (2020).
\newblock Platelet gene expression and function in patients with COVID-19.
\newblock Blood, 136(11), 1317--1329. doi:10.1182/blood.2020007214.

\bibitem[Zaid et al.(2020)]{zaid2020platelets}
Zaid, Y., Puhm, F., Allaeys, I., et al. (2020).
\newblock Platelets can associate with SARS-CoV-2 RNA and are hyperactivated in COVID-19.
\newblock Circulation Research, 127(11), 1404--1418. doi:10.1161/CIRCRESAHA.120.317703.

\bibitem[Hottz et al.(2020)]{hottz2020platelet}
Hottz, E. D., Azevedo-Quintanilha, I. G., Palhinha, L., et al. (2020).
\newblock Platelet activation and platelet-monocyte aggregate formation trigger tissue factor expression in patients with severe COVID-19.
\newblock Blood, 136(11), 1330--1341. doi:10.1182/blood.2020007252.

\bibitem[Middleton et al.(2020)]{middleton2020nets}
Middleton, E. A., He, X.-Y., Denorme, F., et al. (2020).
\newblock Neutrophil extracellular traps contribute to immunothrombosis in COVID-19 acute respiratory distress syndrome.
\newblock Blood, 136(10), 1169--1179. doi:10.1182/blood.2020007008.

\bibitem[Nicolai et al.(2020)]{nicolai2020immunothrombotic}
Nicolai, L., Leunig, A., Brambs, S., et al. (2020).
\newblock Immunothrombotic dysregulation in COVID-19 pneumonia is associated with respiratory failure and coagulopathy.
\newblock Circulation, 142(12), 1176--1189. doi:10.1161/CIRCULATIONAHA.120.048488.

\bibitem[Levi et al.(2020)]{levi2020coagulation}
Levi, M., Thachil, J., Iba, T., \& Levy, J. H. (2020).
\newblock Coagulation abnormalities and thrombosis in patients with COVID-19.
\newblock The Lancet Haematology, 7(6), e438--e440. doi:10.1016/S2352-3026(20)30145-9.

\bibitem[Ackermann et al.(2020)]{ackermann2020pulmonary}
Ackermann, M., Verleden, S. E., Kuehnel, M., et al. (2020).
\newblock Pulmonary vascular endothelialitis, thrombosis, and angiogenesis in Covid-19.
\newblock New England Journal of Medicine, 383(2), 120--128. doi:10.1056/NEJMoa2015432.

\bibitem[Barrett et al.(2021)]{barrett2021platelets}
Barrett, T. J., Bilaloglu, S., Cornwell, M., et al. (2021).
\newblock Platelets contribute to disease severity in COVID-19.
\newblock Journal of Thrombosis and Haemostasis, 19(12), 3139--3153. doi:10.1111/jth.15534.

\bibitem[Qiu et al.(2024)]{qiu2024deciphering}
Qiu, X., Nair, M. G., Jaroszewski, L., \& Godzik, A. (2024).
\newblock Deciphering abnormal platelet subpopulations in COVID-19, sepsis and systemic lupus erythematosus through machine learning and single-cell transcriptomics.
\newblock International Journal of Molecular Sciences, 25(11), 5941. doi:10.3390/ijms25115941.

\bibitem[Adduri et al.(2025)]{adduri2025state}
Adduri, A. K., Gautam, D., Bevilacqua, B., et al. (2025).
\newblock Predicting cellular responses to perturbation across diverse contexts with STATE.
\newblock bioRxiv, 2025.06.26.661135. doi:10.1101/2025.06.26.661135.

\bibitem[Kalfon et al.(2025)]{kalfon2025scprint}
Kalfon, J., Samaran, J., Peyr\'{e}, G., \& Cantini, L. (2025).
\newblock scPRINT: pre-training on 50 million cells allows robust gene network predictions.
\newblock Nature Communications, 16, 3607. doi:10.1038/s41467-025-58699-1.

\bibitem[Pearce et al.(2026)]{pearce2025transcriptformer}
Pearce, J. D., Simmonds, S. E., Mahmoudabadi, G., et al. (2026).
\newblock TranscriptFormer: a generative cell atlas across 1.5 billion years of evolution.
\newblock Science, 393(6806), aec8514. doi:10.1126/science.aec8514.

\bibitem[Van Calster et al.(2019)]{vancalster2019calibration}
Van Calster, B., McLernon, D. J., van Smeden, M., et al. (2019).
\newblock Calibration: the Achilles heel of predictive analytics.
\newblock BMC Medicine, 17(1), 230. doi:10.1186/s12916-019-1466-7.

\bibitem[Collins et al.(2024)]{collins2024tripodai}
Collins, G. S., Moons, K. G. M., Dhiman, P., et al. (2024).
\newblock TRIPOD+AI statement: updated guidance for reporting clinical prediction models that use regression or machine learning methods.
\newblock BMJ, 385, e078378. doi:10.1136/bmj-2023-078378.

\bibitem[Singer et al.(2016)]{singer2016sepsis3}
Singer, M., Deutschman, C. S., Seymour, C. W., et al. (2016).
\newblock The Third International Consensus Definitions for Sepsis and Septic Shock (Sepsis-3).
\newblock JAMA, 315(8), 801--810. doi:10.1001/jama.2016.0287.

\bibitem[Rudd et al.(2020)]{rudd2020global}
Rudd, K. E., Johnson, S. C., Agesa, K. M., et al. (2020).
\newblock Global, regional, and national sepsis incidence and mortality, 1990--2017: analysis for the Global Burden of Disease Study.
\newblock The Lancet, 395(10219), 200--211. doi:10.1016/S0140-6736(19)32989-7.

\bibitem[Reyes et al.(2020)]{reyes2020immune}
Reyes, M., Filbin, M. R., Bhattacharyya, R. P., et al. (2020).
\newblock An immune-cell signature of bacterial sepsis.
\newblock Nature Medicine, 26(3), 333--340. doi:10.1038/s41591-020-0752-4.

\bibitem[Davenport et al.(2016)]{davenport2016genomic}
Davenport, E. E., Burnham, K. L., Radhakrishnan, J., et al. (2016).
\newblock Genomic landscape of the individual host response and outcomes in sepsis: a prospective cohort study.
\newblock The Lancet Respiratory Medicine, 4(4), 259--271. doi:10.1016/S2213-2600(16)00046-1.

\bibitem[Sweeney et al.(2018)]{sweeney2018unsupervised}
Sweeney, T. E., Azad, T. D., Donato, M., et al. (2018).
\newblock Unsupervised analysis of transcriptomics in bacterial sepsis across multiple datasets reveals three robust clusters.
\newblock Critical Care Medicine, 46(6), 915--925. doi:10.1097/CCM.0000000000003084.

\bibitem[Wilk et al.(2020)]{wilk2020singlecell}
Wilk, A. J., Rustagi, A., Zhao, N. Q., et al. (2020).
\newblock A single-cell atlas of the peripheral immune response in patients with severe COVID-19.
\newblock Nature Medicine, 26(7), 1070--1076. doi:10.1038/s41591-020-0944-y.

\bibitem[Schulte-Schrepping et al.(2020)]{schulteschrepping2020severe}
Schulte-Schrepping, J., Reusch, N., Paclik, D., et al. (2020).
\newblock Severe COVID-19 is marked by a dysregulated myeloid cell compartment.
\newblock Cell, 182(6), 1419--1440.e23. doi:10.1016/j.cell.2020.08.001.

\bibitem[Bernardes et al.(2020)]{bernardes2020longitudinal}
Bernardes, J. P., Mishra, N., Tran, F., et al. (2020).
\newblock Longitudinal multi-omics analyses identify responses of megakaryocytes, erythroid cells, and plasmablasts as hallmarks of severe COVID-19.
\newblock Immunity, 53(6), 1296--1314.e9. doi:10.1016/j.immuni.2020.11.017.

\bibitem[Ren et al.(2021)]{ren2021covid}
Ren, X., Wen, W., Fan, X., et al. (2021).
\newblock COVID-19 immune features revealed by a large-scale single-cell transcriptome atlas.
\newblock Cell, 184(7), 1895--1913.e19. doi:10.1016/j.cell.2021.01.053.

\bibitem[Stephenson et al.(2021)]{stephenson2021singlecell}
Stephenson, E., Reynolds, G., Botting, R. A., et al. (2021).
\newblock Single-cell multi-omics analysis of the immune response in COVID-19.
\newblock Nature Medicine, 27(5), 904--916. doi:10.1038/s41591-021-01329-2.

\bibitem[Shappell et al.(2023)]{shappell2023sars}
Shappell, C. N., Klompas, M., Chan, C., et al. (2023).
\newblock Use of electronic clinical data to track incidence and mortality for SARS-CoV-2-associated sepsis.
\newblock JAMA Network Open, 6(9), e2335728. doi:10.1001/jamanetworkopen.2023.35728.

\bibitem[Pedregosa et al.(2011)]{pedregosa2011scikit}
Pedregosa, F., Varoquaux, G., Gramfort, A., et al. (2011).
\newblock Scikit-learn: machine learning in Python.
\newblock Journal of Machine Learning Research, 12, 2825--2830.

\bibitem[Wolf et al.(2018)]{wolf2018scanpy}
Wolf, F. A., Angerer, P., \& Theis, F. J. (2018).
\newblock SCANPY: large-scale single-cell gene expression data analysis.
\newblock Genome Biology, 19(1), 15. doi:10.1186/s13059-017-1382-0.

\bibitem[Luecken \& Theis(2019)]{luecken2019bestpractices}
Luecken, M. D., \& Theis, F. J. (2019).
\newblock Current best practices in single-cell RNA-seq analysis: a tutorial.
\newblock Molecular Systems Biology, 15(6), e8746. doi:10.15252/msb.20188746.

\bibitem[Brier(1950)]{brier1950verification}
Brier, G. W. (1950).
\newblock Verification of forecasts expressed in terms of probability.
\newblock Monthly Weather Review, 78(1), 1--3.

\bibitem[Naeini et al.(2015)]{naeini2015obtaining}
Naeini, M. P., Cooper, G. F., \& Hauskrecht, M. (2015).
\newblock Obtaining well calibrated probabilities using Bayesian binning.
\newblock Proceedings of the AAAI Conference on Artificial Intelligence, 29(1), 2901--2907. doi:10.1609/aaai.v29i1.9602.

\bibitem[Guo et al.(2017)]{guo2017calibration}
Guo, C., Pleiss, G., Sun, Y., \& Weinberger, K. Q. (2017).
\newblock On calibration of modern neural networks.
\newblock In Proceedings of the 34th International Conference on Machine Learning, PMLR 70, 1321--1330.

\bibitem[Niculescu-Mizil \& Caruana(2005)]{niculescu2005predicting}
Niculescu-Mizil, A., \& Caruana, R. (2005).
\newblock Predicting good probabilities with supervised learning.
\newblock In Proceedings of the 22nd International Conference on Machine Learning, 625--632. doi:10.1145/1102351.1102430.

\bibitem[DeLong et al.(1988)]{delong1988comparing}
DeLong, E. R., DeLong, D. M., \& Clarke-Pearson, D. L. (1988).
\newblock Comparing the areas under two or more correlated receiver operating characteristic curves: a nonparametric approach.
\newblock Biometrics, 44(3), 837--845. doi:10.2307/2531595.

\bibitem[Hanley \& McNeil(1982)]{hanley1982meaning}
Hanley, J. A., \& McNeil, B. J. (1982).
\newblock The meaning and use of the area under a receiver operating characteristic (ROC) curve.
\newblock Radiology, 143(1), 29--36. doi:10.1148/radiology.143.1.7063747.

\bibitem[Benjamini \& Hochberg(1995)]{benjamini1995controlling}
Benjamini, Y., \& Hochberg, Y. (1995).
\newblock Controlling the false discovery rate: a practical and powerful approach to multiple testing.
\newblock Journal of the Royal Statistical Society: Series B (Methodological), 57(1), 289--300.

\bibitem[McDermott et al.(2021)]{mcdermott2021reproducibility}
McDermott, M. B. A., Wang, S., Marinsek, N., et al. (2021).
\newblock Reproducibility in machine learning for health research: still a ways to go.
\newblock Science Translational Medicine, 13(586), eabb1655. doi:10.1126/scitranslmed.abb1655.

\bibitem[Luecken et al.(2022)]{luecken2022benchmarking}
Luecken, M. D., B\"{u}ttner, M., Chaichoompu, K., et al. (2022).
\newblock Benchmarking atlas-level data integration in single-cell genomics.
\newblock Nature Methods, 19(1), 41--50. doi:10.1038/s41592-021-01336-8.

\bibitem[Subramanian et al.(2005)]{subramanian2005gsea}
Subramanian, A., Tamayo, P., Mootha, V. K., et al. (2005).
\newblock Gene set enrichment analysis: a knowledge-based approach for interpreting genome-wide expression profiles.
\newblock Proceedings of the National Academy of Sciences, 102(43), 15545--15550. doi:10.1073/pnas.0506580102.

\bibitem[Fang et al.(2023)]{fang2023gseapy}
Fang, Z., Liu, X., \& Peltz, G. (2023).
\newblock GSEApy: a comprehensive package for performing gene set enrichment analysis in Python.
\newblock Bioinformatics, 39(1), btac757. doi:10.1093/bioinformatics/btac757.

\bibitem[Luecken et al.(2025)]{luecken2025openproblems}
Luecken, M. D., Gigante, S., Burkhardt, D. B., et al. (2025).
\newblock Defining and benchmarking open problems in single-cell analysis.
\newblock Nature Biotechnology, 43(7), 1035--1040. doi:10.1038/s41587-025-02694-w.

\bibitem[Corpas et al.(2025)]{corpas2025diversity}
Corpas, M., et al. (2025).
\newblock Bridging genomics' greatest challenge: the diversity gap.
\newblock Cell Genomics, 5, 100724.

\bibitem[Pot et al.(2019)]{pot2019gender}
Pot, M., Spahl, W., \& Prainsack, B. (2019).
\newblock The gender of biomedical data: challenges for personalised and precision medicine.
\newblock Somatechnics, 9, 170--187.

\bibitem[Jiang et al.(2026)]{jiang2026landscape}
Jiang, S., Fang, Z., Zhang, Y., et al. (2026).
\newblock The landscape of single-cell foundation models: design principles, applications, and open challenges.
\newblock Preprints, 2026081166. doi:10.20944/preprints202608.1166.v1.

\bibitem[Bendidi et al.(2024)]{bendidi2024pca}
Bendidi, I., Whitfield, S., Kenyon-Dean, K., et al. (2024).
\newblock Benchmarking transcriptomics foundation models for perturbation analysis: one PCA still rules them all.
\newblock arXiv preprint, 2410.13956. doi:10.48550/arXiv.2410.13956.

\bibitem[DenAdel et al.(2026)]{denadel2026scaling}
DenAdel, A., Hughes, M., Thoutam, A., et al. (2026).
\newblock Evaluating the role of pretraining dataset size and diversity on single-cell foundation model performance.
\newblock Nature Methods, 23, 1447--1457. doi:10.1038/s41592-026-03120-y.

\bibitem[Lee et al.(2020)]{lee2020immunophenotyping}
Lee, J. S., Park, S., Jeong, H. W., et al. (2020).
\newblock Immunophenotyping of COVID-19 and influenza highlights the role of type I interferons in development of severe COVID-19.
\newblock Science Immunology, 5(49), eabd1554. doi:10.1126/sciimmunol.abd1554.

\bibitem[Arunachalam et al.(2020)]{arunachalam2020systems}
Arunachalam, P. S., Wimmers, F., Mok, C. K. P., et al. (2020).
\newblock Systems biological assessment of immunity to mild versus severe COVID-19 infection in humans.
\newblock Science, 369(6508), 1210--1220. doi:10.1126/science.abc6261.

\bibitem[Liu et al.(2021)]{liu2021timeresolved}
Liu, C., Martins, A. J., Lau, W. W., et al. (2021).
\newblock Time-resolved systems immunology reveals a late juncture linked to fatal COVID-19.
\newblock Cell, 184(7), 1836--1857.e22. doi:10.1016/j.cell.2021.02.018.

\bibitem[Combes et al.(2021)]{combes2021global}
Combes, A. J., Courau, T., Kuhn, N. F., et al. (2021).
\newblock Global absence and targeting of protective immune states in severe COVID-19.
\newblock Nature, 591(7848), 124--130. doi:10.1038/s41586-021-03234-7.

\bibitem[Qiu et al.(2021)]{qiu2021sepsis}
Qiu, X., Li, J., Bonenfant, J., et al. (2021).
\newblock Dynamic changes in human single-cell transcriptional signatures during fatal sepsis.
\newblock Journal of Leukocyte Biology, 110(6), 1253--1268. doi:10.1002/JLB.5MA0721-825R.

\bibitem[Jiang et al.(2020)]{jiang2020ards}
Jiang, Y., Rosborough, B. R., Chen, J., et al. (2020).
\newblock Single cell RNA sequencing identifies an early monocyte gene signature in acute respiratory distress syndrome.
\newblock JCI Insight, 5(13), e135678. doi:10.1172/jci.insight.135678.

% Review-content-table studies (P0_onepager.md Sec.5; identifiers verified
% 2026-09-09. Author lists to be completed at the references pass.)
\bibitem[EGSP(2026)]{egsp2026survival}
(2026).
\newblock Leveraging single-cell foundation models for accurate survival outcome prediction.
\newblock Bioinformatics Advances. PMC13032892 (peer-reviewed version of record).

\bibitem[Empirical Evaluation of scFMs(2025)]{empiricalscfm2025cancer}
(2025).
\newblock Empirical Evaluation of Single-Cell Foundation Models for Predicting Cancer Outcomes.
\newblock bioRxiv preprint, 2025.10.31.685892. doi:10.1101/2025.10.31.685892.

\bibitem[TissueFormer(2025)]{tissueformer2025covid}
(2025).
\newblock TissueFormer: COVID-19 severity classification with a murine-adapted Geneformer.
\newblock bioRxiv preprint, 2025.08.17.670735. doi:10.1101/2025.08.17.670735.

\end{thebibliography}

\end{document}
